%!TEX root = forallxyyc.tex
\part{Tabelas-verdade}
\label{ch.TruthTables}
\addtocontents{toc}{\protect\mbox{}\protect\hrulefill\par}

\chapter{Tabelas-verdade características}
\label{s:CharacteristicTruthTables}

Qualquer sentença da LVF é composta de letras sentenciais, possivelmente combinadas por meio de conectivos sentenciais. O valor de verdade da sentença composta depende somente dos valores de verdade das letras sentenciais que a compõem. Para saber o valor de verdade de \textquotedblleft$(D \eand E)$\textquotedblright{}, por exemplo, você só precisa saber o valor de verdade de \textquotedblleft$D$\textquotedblright{} e o valor de verdade de \textquotedblleft$E$\textquotedblright{}.

Introduzimos cinco conectivos em \cref{s:TFLConnectives}. Portanto, só precisamos explicar como eles mapeiam valores de verdade. Por conveniência, abreviaremos \textquotedblleft{}Verdadeiro\textquotedblright{} por \textquotedblleft T\textquotedblright{} e \textquotedblleft{}Falso\textquotedblright{} por \textquotedblleft F\textquotedblright{}. (Mas, para deixar claro, os dois valores de verdade são Verdadeiro e Falso; os valores de verdade não são \emph{letras}!)

\newglossaryentry{truth value}
                 {
                   name = valor de verdade,
                   description = {Um dos dois valores lógicos que as sentenças podem ter: Verdadeiro e Falso}
                   }

\paragraph{Negação} Para qualquer sentença \metav{A}: se \metav{A} é verdadeira, então \enot\metav{A} é falsa; e se \enot\metav{A} é verdadeira, então \metav{A} é falsa. Podemos resumir isso na \emph{tabela-verdade característica} da negação:
\begin{center}
\begin{tabular}{c|c}
\metav{A} & \enot\metav{A}\\
\hline
T & F\\
F & T 
\end{tabular}
\end{center}

\paragraph{Conjunção} Para quaisquer sentenças \metav{A} e \metav{B}, \metav{A}\eand\metav{B} é verdadeira se e somente se tanto \metav{A} quanto \metav{B} são verdadeiras. Podemos resumir isso na {tabela-verdade característica} da conjunção:
\begin{center}
\begin{tabular}{c c |c}
\metav{A} & \metav{B} & $\metav{A}\eand\metav{B}$\\
\hline
T & T & T\\
T & F & F\\
F & T & F\\
F & F & F
\end{tabular}
\end{center}
Observe que o valor de verdade de $\metav{A} \eand \metav{B}$ é sempre
o mesmo que o valor de verdade de $\metav{B} \eand \metav{A}$.
Conectivos que têm essa propriedade são chamados de \emph{comutativos}.

\paragraph{Disjunção} Recorde que \textquotedblleft$\eor$\textquotedblright{} sempre representa ou inclusivo. Portanto, para quaisquer sentenças \metav{A} e \metav{B}, $\metav{A}\eor \metav{B}$ é verdadeira se e somente se \metav{A} ou \metav{B} é verdadeira. Podemos resumir isso na {tabela-verdade característica} da disjunção:
\begin{center}
\begin{tabular}{c c|c}
\metav{A} & \metav{B} & $\metav{A}\eor\metav{B}$ \\
\hline
T & T & T\\
T & F & T\\
F & T & T\\
F & F & F
\end{tabular}
\end{center}
Assim como a conjunção, a disjunção é comutativa.

\paragraph{Condicional} Vamos simplesmente admitir:
Os condicionais são uma bagunça. Exatamente qual é o grau dessa
bagunça é uma questão \emph{filosoficamente} controversa. Discutiremos
algumas das sutilezas em
\cref{s:IndicativeSubjunctive,s:ParadoxesOfMaterialConditional}. Por
ora, estipularemos o seguinte: $\metav{A}\eif\metav{B}$ é falsa se e somente
se \metav{A} é verdadeira e \metav{B} é falsa. O conectivo com essas
condições de verdade estipuladas é chamado de
\define{condicional material}. Podemos resumir isso com uma tabela-verdade
característica do condicional.
\begin{center}
\begin{tabular}{c c|c}
\metav{A} & \metav{B} & $\metav{A}\eif\metav{B}$\\
\hline
T & T & T\\
T & F & F\\
F & T & T\\
F & F & T
\end{tabular}
\end{center}
O condicional \emph{não} é comutativo. Não se pode trocar o antecedente
e o consequente sem mudar o significado da sentença; $\metav{A}\eif\metav{B}$
e $\metav{B} \eif \metav{A}$ têm tabelas-verdade diferentes.

\paragraph{Bicondicional} Como um bicondicional deve ser o mesmo que a conjunção dos condicionais que percorrem ambas as direções, queremos que a tabela-verdade do bicondicional seja:
\begin{center}
\begin{tabular}{c c|c}
\metav{A} & \metav{B} & $\metav{A}\eiff\metav{B}$\\
\hline
T & T & T\\
T & F & F\\
F & T & F\\
F & F & T
\end{tabular}
\end{center}
O conectivo com esta tabela-verdade também é frequentemente chamado de
\define{bicondicional material}. Como era de esperar, o bicondicional é
comutativo.
\chapter{Conectivos verofuncionais}
\label{s:TruthFunctionality}

\section{A ideia da verofuncionalidade}
Introduzamos uma ideia importante.
	\factoidbox{
		Um conectivo é \define{verofuncional} \ifeff{} o valor de verdade de uma sentença que tem esse conectivo como seu operador lógico principal é determinado univocamente pelo(s) valor(es) de verdade da(s) sentença(s) constituintes.
	}
\newglossaryentry{truth-functional connective}
{
name=conectivo verofuncional,
description={Um operador que constrói sentenças maiores a partir de sentenças menores e fixa o \gls{truth value} da sentença resultante com base apenas no valor de verdade das sentenças componentes}
}
        
Todo conectivo da LVF é verofuncional. O valor de verdade de uma negação é determinado univocamente pelo valor de verdade da sentença não negada. O valor de verdade de uma conjunção é determinado univocamente pelo valor de verdade de ambos os conjuntivos. O valor de verdade de uma disjunção é determinado univocamente pelo valor de verdade de ambos os disjuntivos, e assim por diante. Para determinar o valor de verdade de uma sentença da LVF, só precisamos conhecer o valor de verdade de seus componentes.

É isso que dá nome à LVF: ela é uma \emph{lógica proposicional verofuncional}.

Muitas línguas usam conectivos que não são verofuncionais. Em português, por exemplo, podemos formar uma nova sentença a partir de qualquer sentença mais simples prefixando-a com \textquotedblleft{}É necessariamente o caso que\ldots\textquotedblright{}. O valor de verdade dessa nova sentença não é fixado apenas pelo valor de verdade da sentença original. Considere duas sentenças verdadeiras:
	\begin{compactlist}
		\item $2 + 2 = 4$
		\item Shostakovich compôs quinze quartetos de cordas.
	\end{compactlist}
Embora seja necessariamente o caso que $2 + 2 = 4$, não é \emph{necessariamente} o caso que Shostakovich compôs quinze quartetos de cordas. Se Shostakovich tivesse morrido antes, não teria conseguido terminar o Quarteto nº~15; se tivesse vivido mais, talvez tivesse composto mais alguns. Portanto, \textquotedblleft{}É necessariamente o caso que\ldots\textquotedblright{} não é \emph{verofuncional}.

\section{Simbolizar versus traduzir}
Todos os conectivos da LVF são verofuncionais, mas, mais do que isso: eles realmente não fazem \emph{nada além de} nos mapear entre valores de verdade.

Ao simbolizarmos uma sentença ou um argumento na LVF, ignoramos tudo \emph{além} da contribuição que os valores de verdade de um componente podem dar ao valor de verdade do todo. Há sutilezas em nossas afirmações ordinárias que vão muito além de seus meros valores de verdade. Sarcasmo; poesia; implicatura maliciosa; ênfase; tudo isso são partes importantes do discurso cotidiano, mas nada disso é preservado na LVF. Como observamos em \cref{s:TFLConnectives}, a LVF não consegue captar as diferenças sutis entre as seguintes sentenças em português:
	\begin{compactlist}
		\item Dana é uma lógica e Dana é uma pessoa gentil
		\item Embora Dana seja uma lógica, Dana é uma pessoa gentil
		\item Dana é uma lógica apesar de ser uma pessoa gentil
		\item Dana é uma pessoa gentil, mas também é uma lógica
		\item Não obstante Dana ser uma lógica, ela é uma pessoa gentil
	\end{compactlist}
Todas as sentenças acima serão simbolizadas com a mesma sentença da LVF, talvez \textquotedblleft{}$L \eand N$\textquotedblright{}.

Continuamos dizendo que usamos sentenças da LVF para \emph{simbolizar} sentenças em português. Muitos outros livros didáticos falam em \emph{traduzir} sentenças em português para a LVF. Entretanto, uma boa tradução deveria preservar certos aspectos do significado e — como acabamos de ver — a LVF simplesmente não consegue fazer isso. É por isso que falaremos em \emph{simbolizar} sentenças em português, em vez de \emph{traduzi-las}.

Isso afeta a maneira como devemos entender nossas chaves de simbolização. Considere uma chave como:
	\begin{ekey}
		\item[L] Dana é uma lógica.
		\item[N] Dana é uma pessoa gentil.
	\end{ekey}
Outros livros didáticos entenderão isso como uma estipulação de que a sentença da LVF \textquotedblleft{}$L$\textquotedblright{} deve \emph{significar} que Dana é uma lógica e de que a sentença da LVF \textquotedblleft{}$N$\textquotedblright{} deve \emph{significar} que Dana é uma pessoa gentil. Mas a LVF é totalmente incapaz de lidar com \emph{significado}. A chave de simbolização anterior não faz nada mais nem menos do que estipular que a sentença da LVF \textquotedblleft{}$L$\textquotedblright{} deve ter o mesmo valor de verdade que a sentença em português \textquotedblleft{}Dana é uma lógica\textquotedblright{} (seja lá qual for esse valor), e que a sentença da LVF \textquotedblleft{}$N$\textquotedblright{} deve ter o mesmo valor de verdade que a sentença em português \textquotedblleft{}Dana é uma pessoa gentil\textquotedblright{} (seja lá qual for esse valor).
	\factoidbox{
		Quando tratamos uma sentença da LVF como \emph{simbolizando} uma sentença em português, estamos estipulando que a sentença da LVF deve ter o mesmo valor de verdade que essa sentença em português.
	}


\section{Condicionais indicativos versus subjuntivos}\label{s:IndicativeSubjunctive}
Queremos enfatizar o ponto de que a LVF só pode lidar com funções de verdade considerando o caso do condicional. Quando introduzimos a tabela-verdade característica do condicional material em \cref{s:CharacteristicTruthTables}, não dissemos nada para justificá-la. Ofereçamos agora uma justificativa, seguindo Dorothy Edgington.\footnote{Dorothy Edgington, \textquotedblleft{}Conditionals\textquotedblright{},
\textit{Stanford Encyclopedia of Philosophy} (Fall 2020)
(\url{https://plato.stanford.edu/archives/fall2020/entries/conditionals/}).}

Suponha que Lara tenha desenhado algumas formas em uma folha de papel e colorido algumas delas. Não as vimos, mas ainda assim afirmamos:
	\begin{quote}
		Se alguma forma é cinza, então essa forma também é circular.
	\end{quote}
Acontece que Lara desenhou o seguinte:
\begin{center}
\begin{tikzpicture}[alt={Há três formas: a primeira forma é um círculo cinza rotulado A, há algum espaço vazio, a segunda forma é um círculo branco rotulado C e a quarta é um quadrado branco rotulado D.}]
	\node[circle, grey_shape] (cat1) {A};
	\node[right=10pt of cat1, diamond, phantom_shape] (cat2)  { } ;
	\node[right=10pt of cat2, circle, white_shape] (cat3)  {C} ;
	\node[right=10pt of cat3, white_shape] (cat4)  {D};
\end{tikzpicture}
\bmlDescription{Há três formas: a primeira forma é um círculo cinza rotulado A, há algum espaço vazio, a segunda forma é um círculo branco rotulado C e a quarta é um quadrado branco rotulado D.}
\end{center}
Neste caso, nossa afirmação é certamente verdadeira. As formas C e D não são cinzas e, portanto, dificilmente podem constituir \emph{contraexemplos} à nossa afirmação. A forma A \emph{é} cinza, mas felizmente também é circular. Portanto, nossa afirmação não tem contraexemplos. Ela deve ser verdadeira. Isso significa que cada uma das seguintes \emph{instâncias} da nossa afirmação também deve ser verdadeira:
	\begin{itemize}
		\item Se A é cinza, então é circular

		\qquad (antecedente verdadeiro, consequente verdadeiro)
		\item Se C é cinza, então é circular

		\qquad (antecedente falso, consequente verdadeiro)
		\item Se D é cinza, então é circular

		\qquad (antecedente falso, consequente falso)
	\end{itemize}
Contudo, se Lara tivesse desenhado uma quarta forma, assim:
\begin{center}
\begin{tikzpicture}[alt={Há quatro formas: a primeira forma é um círculo cinza rotulado A, a segunda é um quadrado cinza rotulado B, a terceira é um círculo branco rotulado C e a quarta é um quadrado branco rotulado D.}]
	\node[circle, grey_shape] (cat1) {A};
	\node[right=10pt of cat1, grey_shape] (cat2)  {B};
	\node[right=10pt of cat2, circle, white_shape] (cat3)  {C};
	\node[right=10pt of cat3, white_shape] (cat4)  {D};
\end{tikzpicture}
\bmlDescription{Há quatro formas: a primeira forma é um círculo cinza rotulado A, a segunda é um quadrado cinza rotulado B, a terceira é um círculo branco rotulado C e a quarta é um quadrado branco rotulado D.}
\end{center}
então nossa afirmação teria sido falsa. Portanto, esta instância da nossa afirmação também deve ser falsa:
	\begin{itemize}
		\item Se B é cinza, então é circular

		\qquad (antecedente verdadeiro, consequente falso)
	\end{itemize}
Agora, lembre-se de que todo conectivo da LVF precisa ser verofuncional. Isso significa que os valores de verdade do antecedente e do consequente, sozinhos, devem determinar univocamente o valor de verdade do condicional como um todo. Assim, a partir dos valores de verdade das nossas quatro afirmações — que nos fornecem todas as combinações possíveis de verdade e falsidade no antecedente e no consequente — podemos deduzir a tabela-verdade do condicional material.

O que esse argumento mostra é que \textquotedblleft{}$\eif$\textquotedblright{} é o \emph{melhor} candidato a um condicional verofuncional. Dito de outro modo, \emph{é o melhor condicional que a LVF pode oferecer}. Mas ele é bom como substituto dos condicionais que usamos na linguagem cotidiana? Considere duas sentenças:
	\begin{enumerate}
		\item\label{brownwins1} Se Hillary Clinton tivesse vencido a eleição de 2016, então teria sido a primeira mulher presidente dos EUA.
		\item\label{brownwins2} Se Hillary Clinton tivesse vencido a eleição de 2016, então teria se transformado em um balão cheio de hélio e flutuado para longe no céu noturno.
	\end{enumerate}
\Cref*{brownwins1} é verdadeira; \cref*{brownwins2} é falsa, mas ambas têm antecedentes falsos e consequentes falsos. (Hillary não venceu; não se tornou a primeira mulher presidente dos EUA; e não se encheu de hélio nem saiu flutuando.) Portanto, o valor de verdade da sentença inteira não é determinado univocamente pelo valor de verdade de suas partes.

O ponto crucial é que \cref*{brownwins1,brownwins2} empregam condicionais \emph{subjuntivos}, em vez de condicionais \emph{indicativos}. Eles nos pedem que imaginemos algo contrário aos fatos — afinal, Hillary Clinton perdeu a eleição de 2016 — e então nos pedem que avaliemos o que \emph{teria} acontecido nesse caso. Tais considerações simplesmente não podem ser tratadas usando~\textquotedblleft{}$\eif$\textquotedblright{}.

Diremos mais sobre as dificuldades com os condicionais em \cref{s:ParadoxesOfMaterialConditional}. Por ora, vamos nos contentar com a observação de que \textquotedblleft{}$\eif$\textquotedblright{} é o único candidato a condicional verofuncional para a LVF, mas que muitos condicionais em português não podem ser representados adequadamente usando \textquotedblleft{}$\eif$\textquotedblright{}. A LVF é uma linguagem intrinsecamente limitada.

\chapter{Tabelas-verdade completas}
\label{s:CompleteTruthTables}

Até aqui, usamos chaves de simbolização para atribuir valores de verdade a sentenças da LVF \emph{indiretamente}. Por exemplo, podemos dizer que a sentença da LVF \textquotedblleft{}$B$\textquotedblright{} deve ser verdadeira \ifeff{} Big Ben está em Londres. Como Big Ben \emph{está} em Londres, essa simbolização tornaria \textquotedblleft{}$B$\textquotedblright{} verdadeira. Mas também podemos atribuir valores de verdade \emph{diretamente}. Podemos simplesmente estipular que \textquotedblleft{}$B$\textquotedblright{} é verdadeira ou estipular que ela é falsa. Essas estipulações são chamadas de \emph{valorações}:
	\factoidbox{
		Uma \define{valora\c{c}\~{a}o} é qualquer atribuição de valores de verdade a determinadas letras sentenciais da LVF.
	}

\newglossaryentry{valuation}
{
name=valoração,
description={Uma atribuição de \glspl{truth value} a determinadas \glspl{sentence letter}}
}

A força das tabelas-verdade está no seguinte. Cada linha de uma tabela-verdade representa uma valoração possível. A tabela-verdade completa representa todas as valorações possíveis. E a tabela-verdade nos fornece um meio de calcular o valor de verdade de sentenças complexas, para cada valoração possível. Mas tudo isso é mais fácil de explicar com um exemplo.

\section{Um exemplo trabalhado}
Considere a sentença \textquotedblleft{}$(H\eand I)\eif H$\textquotedblright{}. Há quatro maneiras possíveis de atribuir Verdadeiro e Falso à letra sentencial \textquotedblleft{}$H$\textquotedblright{} e à letra \textquotedblleft{}$I$\textquotedblright{} — quatro valorações — que podemos representar assim:
\begin{center}
\begin{tabular}{c c|d e e e f}
$H$&$I$&$(H$&\eand&$I)$&\eif&$H$\\
\hline
 T & T\\
 T & F\\
 F & T\\
 F & F
\end{tabular}
\end{center}
Para calcular o valor de verdade da sentença inteira \textquotedblleft{}$(H \eand I) \eif H$\textquotedblright{}, primeiro copiamos os valores de verdade das letras sentenciais e os escrevemos abaixo das letras na sentença:
\begin{center}
\begin{tabular}{c c|d e e e f}
$H$&$I$&$(H$&\eand&$I)$&\eif&$H$\\
\hline
 T & T & {T} & & {T} & & {T}\\
 T & F & {T} & & {F} & & {T}\\
 F & T & {F} & & {T} & & {F}\\
 F & F & {F} & & {F} & & {F}
\end{tabular}
\end{center}
Agora considere a subsentença \textquotedblleft{}$(H\eand I)$\textquotedblright{}. Ela é uma conjunção, $(\metav{A}\eand\metav{B})$, com \textquotedblleft{}$H$\textquotedblright{} como \metav{A} e \textquotedblleft{}$I$\textquotedblright{} como \metav{B}. A tabela-verdade característica da conjunção fornece as condições de verdade para \emph{qualquer} sentença da forma $(\metav{A}\eand\metav{B})$, quaisquer que sejam $\metav{A}$ e $\metav{B}$. Ela representa o fato de que uma conjunção é verdadeira \ifeff{} ambos os conjuntivos são verdadeiros. Neste caso, nossos conjuntivos são simplesmente \textquotedblleft{}$H$\textquotedblright{} e \textquotedblleft{}$I$\textquotedblright{}. Ambos são verdadeiros na primeira linha da tabela-verdade, e somente nela. Consequentemente, podemos calcular o valor de verdade da conjunção em todas as quatro linhas.
\begin{center}
\begin{tabular}{c c|d e e e f}
 & & \metav{A} & \eand & \metav{B} & & \\
$H$&$I$&$(H$&\eand&$I)$&\eif&$H$\\
\hline
 T & T & T & {T} & T & & T\\
 T & F & T & {F} & F & & T\\
 F & T & F & {F} & T & & F\\
 F & F & F & {F} & F & & F
\end{tabular}
\end{center}
Agora, a sentença inteira com que estamos lidando é um condicional, $\metav{A}\eif\metav{B}$, com \textquotedblleft{}$(H \eand I)$\textquotedblright{} como \metav{A} e \textquotedblleft{}$H$\textquotedblright{} como \metav{B}. Na segunda linha, por exemplo, \textquotedblleft{}$(H\eand I)$\textquotedblright{} é falsa e \textquotedblleft{}$H$\textquotedblright{} é verdadeira. Como um condicional é verdadeiro quando o antecedente é falso, escrevemos um \textquotedblleft{}T\textquotedblright{} na segunda linha, abaixo do símbolo condicional. Continuamos com as outras três linhas e obtemos:
\begin{center}
\begin{tabular}{c c| d e e e f}
 & &  & \metav{A} &  &\eif &\metav{B} \\
$H$&$I$&$(H$&\eand&$I)$&\eif&$H$\\
\hline
 T & T &  & {T} &  &{T} & T\\
 T & F &  & {F} &  &{T} & T\\
 F & T &  & {F} &  &{T} & F\\
 F & F &  & {F} &  &{T} & F
\end{tabular}
\end{center}
O condicional é o operador lógico principal da sentença, portanto a coluna de \textquotedblleft{}T\textquotedblright{} abaixo do condicional nos diz que a sentença \textquotedblleft{}$(H \eand I)\eif H$\textquotedblright{} é verdadeira independentemente dos valores de verdade de \textquotedblleft{}$H$\textquotedblright{} e \textquotedblleft{}$I$\textquotedblright{}. Eles podem ser verdadeiros ou falsos em qualquer combinação, e a sentença composta ainda resulta verdadeira. Como consideramos todas as quatro atribuições possíveis de verdade e falsidade a \textquotedblleft{}$H$\textquotedblright{} e \textquotedblleft{}$I$\textquotedblright{} — isto é, como consideramos todas as diferentes \emph{valorações} — podemos dizer que \textquotedblleft{}$(H \eand I)\eif H$\textquotedblright{} é verdadeira em toda valoração.

Neste exemplo, não repetimos todas as entradas em cada coluna de cada tabela sucessiva. Ao escrever tabelas-verdade no papel, porém, é impraticável apagar colunas inteiras ou reescrever a tabela toda a cada passo. Embora fique mais carregada, a tabela-verdade pode ser escrita assim:
\begin{center}
\begin{tabular}{c c| d e e e f}
$H$&$I$&$(H$&\eand&$I)$&\eif&$H$\\
\hline
 T & T & T & {T} & T & \TTbf{T} & T\\
 T & F & T & {F} & F & \TTbf{T} & T\\
 F & T & F & {F} & T & \TTbf{T} & F\\
 F & F & F & {F} & F & \TTbf{T} & F
\end{tabular}
\end{center}
A maior parte das colunas abaixo da sentença serve apenas para fins de registro. A coluna que mais importa é a coluna abaixo do \emph{operador lógico principal} da sentença, pois ela informa o valor de verdade da sentença inteira. Destacamos isso colocando essa coluna em negrito. Ao trabalhar com tabelas-verdade por conta própria, você deve destacá-la da mesma maneira (talvez realçando-a).

\section{Construindo tabelas-verdade completas}
Uma \define{tabela-verdade completa} tem uma linha para cada atribuição possível de Verdadeiro e Falso às letras sentenciais relevantes. Cada linha representa uma \emph{valoração}, e uma tabela-verdade completa tem uma linha para todas as diferentes valorações.

\newglossaryentry{complete truth table}
{
name=tabela-verdade completa,
description={Uma tabela que fornece todos os possíveis \glspl{truth value} para uma \gls{sentence of TFL} ou sentenças da LVF, com uma linha para cada \gls{valuation} possível de todas as letras sentenciais}
}

O tamanho da tabela-verdade completa depende do número de letras sentenciais diferentes na tabela. Uma sentença que contém apenas uma letra sentencial requer somente duas linhas, como na tabela-verdade característica da negação. Isso é verdade mesmo que a mesma letra se repita muitas vezes, como na sentença
\textquotedblleft{}$[(C\eiff C) \eif C] \eand \enot(C \eif C)$\textquotedblright{}.
A tabela-verdade completa requer apenas duas linhas porque há somente duas possibilidades: \textquotedblleft{}$C$\textquotedblright{} pode ser verdadeira ou pode ser falsa. A tabela-verdade para essa sentença é assim:
\begin{center}
\begin{tabular}{c| d e e e e e e e e e e e e e e f}
$C$&$[($&$C$&\eiff&$C$&$)$&\eif&$C$&$]$&\eand&\enot&$($&$C$&\eif&$C$&$)$\\
\hline
 T &    & T &  T  & T &   & T  & T & &\TTbf{F}&  F& &   T &  T  & T &   \\
 F &    & F &  T  & F &   & F  & F & &\TTbf{F}&  F& &   F &  T  & F &   \\
\end{tabular}
\end{center}
Ao observar a coluna abaixo do operador lógico principal, vemos que a sentença é falsa nas duas linhas da tabela; isto é, a sentença é falsa independentemente de \textquotedblleft{}$C$\textquotedblright{} ser verdadeira ou falsa. Ela é falsa em toda valoração.

Haverá quatro linhas na tabela-verdade completa para uma sentença contendo duas letras sentenciais, como nas tabelas-verdade características ou na tabela-verdade para \textquotedblleft{}$(H \eand I)\eif H$\textquotedblright{}.

Haverá oito linhas na tabela-verdade completa para uma sentença contendo três letras sentenciais, por exemplo:
\begin{center}
\begin{tabular}{c c c|d e e e f}
$M$&$N$&$P$&$M$&\eand&$(N$&\eor&$P)$\\
\hline
%           M        &     N   v   P
T & T & T & T & \TTbf{T} & T & T & T\\
T & T & F & T & \TTbf{T} & T & T & F\\
T & F & T & T & \TTbf{T} & F & T & T\\
T & F & F & T & \TTbf{F} & F & F & F\\
F & T & T & F & \TTbf{F} & T & T & T\\
F & T & F & F & \TTbf{F} & T & T & F\\
F & F & T & F & \TTbf{F} & F & T & T\\
F & F & F & F & \TTbf{F} & F & F & F
\end{tabular}
\end{center}
A partir dessa tabela, sabemos que a sentença \textquotedblleft{}$M\eand(N\eor P)$\textquotedblright{} pode ser verdadeira ou falsa, dependendo dos valores de verdade de \textquotedblleft{}$M$\textquotedblright{}, \textquotedblleft{}$N$\textquotedblright{} e \textquotedblleft{}$P$\textquotedblright{}.

Uma tabela-verdade completa para uma sentença que contém quatro letras sentenciais diferentes requer 16 linhas. Cinco letras, 32 linhas. Seis letras, 64 linhas. E assim por diante. Para sermos perfeitamente gerais: se uma tabela-verdade completa tem $n$ letras sentenciais diferentes, então ela deve ter $2^n$ linhas.

Para preencher as colunas de uma tabela-verdade completa, comece com a letra sentencial mais à direita e alterne entre \textquotedblleft{}T\textquotedblright{} e \textquotedblleft{}F\textquotedblright{}. Na coluna seguinte à esquerda, escreva dois \textquotedblleft{}T\textquotedblright{}, escreva dois \textquotedblleft{}F\textquotedblright{} e repita. Para a terceira letra sentencial, escreva quatro \textquotedblleft{}T\textquotedblright{} seguidos de quatro \textquotedblleft{}F\textquotedblright{}. Isso produz uma tabela-verdade com oito linhas como a acima. Para uma tabela-verdade com 16 linhas, a coluna seguinte de letras sentenciais deve ter oito \textquotedblleft{}T\textquotedblright{} seguidos de oito \textquotedblleft{}F\textquotedblright{}. Para uma tabela com 32 linhas, a coluna seguinte teria 16 \textquotedblleft{}T\textquotedblright{} seguidos de 16 \textquotedblleft{}F\textquotedblright{} e assim por diante.

\section{Mais sobre parênteses}\label{s:MoreBracketingConventions}
Considere estas duas sentenças:
	\begin{align*}
		((A \eand B) \eand C)\\
		(A \eand (B \eand C))
	\end{align*}
Elas são equivalentes verofuncionalmente. Consequentemente, nunca fará diferença, do ponto de vista do valor de verdade — que é tudo com que a LVF se preocupa (ver \cref{s:TruthFunctionality}) —, qual das duas sentenças afirmamos (ou negamos). Embora a ordem dos parênteses não importe para a verdade delas, não devemos simplesmente eliminá-los. A expressão
	\begin{align*}
		A \eand B \eand C
	\end{align*}
é ambígua entre as duas sentenças acima. A mesma observação vale para disjunções. As sentenças seguintes são logicamente equivalentes:
	\begin{align*}
		((A \eor B) \eor C)\\
		(A \eor (B \eor C))
	\end{align*}
Mas não devemos simplesmente escrever:
	\begin{align*}
		A \eor B \eor C
	\end{align*}
Na verdade, é um fato específico da tabela-verdade característica de
$\eor$ e $\eand$ que garante que quaisquer duas conjunções (ou
disjunções) das mesmas sentenças são equivalentes verofuncionalmente,
independentemente de como você coloca os parênteses. \emph{Isso só é
verdade para conjunções, disjunções e bicondicionais}, contudo. As duas
sentenças seguintes têm tabelas-verdade \emph{diferentes}:
	\begin{align*}
		((A \eif B) \eif C)\\
		(A \eif (B \eif C))
	\end{align*}
Portanto, se escrevêssemos:
	\begin{align*}
		A \eif B \eif C
	\end{align*}
isso seria perigosamente ambíguo. Omitir os parênteses neste caso seria desastroso. Da mesma forma, estas sentenças têm tabelas-verdade diferentes:
	\begin{align*}
		((A \eor B) \eand C)\\
		(A \eor (B \eand C))
	\end{align*}
Portanto, se escrevêssemos:
	\begin{align*}
		A \eor B \eand C
	\end{align*}
isso seria perigosamente ambíguo. \emph{Nunca escreva isso.} A moral é: nunca elimine os parênteses (exceto os mais externos).

\practiceproblems\label{pr.TT.TTorC}
\problempart
Forneça tabelas-verdade completas para cada uma das seguintes sentenças:
\begin{compactlist}
\item $A \eif A$ %taut
\item $C \eif\enot C$ %contingent
\item $(A \eiff B) \eiff \enot(A\eiff \enot B)$ %tautology
\item $(A \eif B) \eor (B \eif A)$ % taut
\item $(A \eand B) \eif (B \eor A)$  %taut
\item $\enot(A \eor B) \eiff (\enot A \eand \enot B)$ %taut
\item $\bigl[(A\eand B) \eand\enot(A\eand B)\bigr] \eand C$ %contradiction
\item $[(A \eand B) \eand C] \eif B$ %taut
\item $\enot\bigl[(C\eor A) \eor B\bigr]$ %contingent
\end{compactlist}
\problempart
Verifique todas as afirmações feitas em \cref{s:MoreBracketingConventions}, isto é, mostre que:
\begin{compactlist}
	\item \textquotedblleft{}$((A \eand B) \eand C)$\textquotedblright{} e \textquotedblleft{}$(A \eand (B \eand C))$\textquotedblright{} têm a mesma tabela-verdade
	\item \textquotedblleft{}$((A \eor B) \eor C)$\textquotedblright{} e \textquotedblleft{}$(A \eor (B \eor C))$\textquotedblright{} têm a mesma tabela-verdade
	\item \textquotedblleft{}$((A \eor B) \eand C)$\textquotedblright{} e \textquotedblleft{}$(A \eor (B \eand C))$\textquotedblright{} não têm a mesma tabela-verdade
	\item \textquotedblleft{}$((A \eif B) \eif C)$\textquotedblright{} e \textquotedblleft{}$(A \eif (B \eif C))$\textquotedblright{} não têm a mesma tabela-verdade
\end{compactlist}
Além disso, verifique se:
\begin{compactlist}
	\item[5.] \textquotedblleft{}$((A \eiff B) \eiff C)$\textquotedblright{} e \textquotedblleft{}$(A \eiff (B \eiff C))$\textquotedblright{} têm a mesma tabela-verdade
\end{compactlist}

\problempart
Escreva tabelas-verdade completas para as sentenças seguintes e marque a coluna que representa os possíveis valores de verdade da sentença inteira.

\begin{compactlist}

\item $\enot (S \eiff (P \eif S))$

%\begin{tabular}{c|c|ccccc}
%\cline{2-2}
%1.	&	\enot 	&	(S 	&	\eiff	&	(P 	&	\eif	&	S))	\\ 
%\cline{2-7}
%	& 	F 		&	T	&	T	&	T	&	T	&	T	\\
%	& 	F 		&	T	&	T	&	F	&	T	&	T	\\
%	& 	F 		&	F	&	T	&	T	&	F	&	F	\\
%	& 	T 		&	F	&	F	&	F	&	T	&	F	\\
%\cline{2-2}
%\end{tabular}


 \item $\enot [(X \eand Y) \eor (X \eor Y)]$

%\begin{tabular}{c|c|ccccccc}
%\cline{2-2}
%2.	&	\enot	&	 [(X 	&	\eand& 	Y) 	&	\eor 	&	(X 	&	\eor 	&	Y)] \\
%\cline{2-9}
%	&	F	&	T	&	T	&	T	&	T	&	T	&	T	&	T	\\
%	&	F	&	T	&	F	&	F	&	T	&	T	&	T	&	F	\\
%	&	F	&	F	&	F	&	T	&	T	&	F	&	T	&	T	\\
%	&	T	&	F	&	F	&	F	&	F	&	F	&	F	&	F	\\
%\cline{2-9}
%\end{tabular}


\item $(A \eif B) \eiff (\enot B\eiff \enot A)$
%\begin{tabular}{cccc|c|ccccc}
%\cline{5-5}
%3.	&	(A 	&	\eif	&	B)	&	 \eiff 	&	(\enot&	B 	&	\eiff 	&	 \enot 	& 	 A) \\
%\cline{2-10}
%	&	T	&	T	&	T	&	T		&	F	 &	T	&	T	&	F		&	T	\\	
%	&	T	&	F	&	F	&	T		&	T	 &	F	&	F	&	F		&	T	\\
%	&	F	&	T	&	T	&	F		&	F	 &	T	&	F	&	T		&	F	\\
%	&	F		&	T		&	F		&	T		&	T	 &	F	&	T	&	T	&	F	\\
%\cline{5-5}
%\end{tabular}

\item $[C \eiff (D \eor E)] \eand \enot C$

%\begin{tabular}{cccccc|c|cc}
%\cline{7-7}
%4.	&	[C 	&	\eiff 	&	(D 	&	\eor 	&	E)] 	&	\eand 	&	 \enot 	&	 C \\
%\cline{2-9}
%	&	T	&	T	&	T	&	T	&	T	&	F		&	F		&	T	\\
%	&	T	&	T	&	T	&	T	&	F	&	F		&	F		&	T	\\
%	&	T	&	T	&	F	&	T	&	T	&	F		&	F		&	T	\\
%	&	T	&	F	&	F	&	F	&	F	&	F		&	F		&	T	\\
%	&	F	&	F	&	T	&	T	&	T	&	F		&	T		&	F	\\
%	&	F	&	F	&	T	&	T	&	F	&	F		&	T		&	F	\\
%	&	F	&	F	&	F	&	T	&	T	&	F		&	T		&	F	\\
%	&	F	&	T	&	F	&	F	&	F	&	T		&	T		&	F	\\
%\cline{7-7}
%\end{tabular}

\item $\enot(G \eand (B \eand H)) \eiff (G \eor (B \eor H))$
%
%\begin{tabular}{ccccccc|c|ccccc}
%\cline{8-8}
%5.	&\enot&	(G 	&\eand &	(B 	&	 \eand 	&	 H))	&	\eiff 	&	(G 	& \eor 	& (B 	& \eor	& H))	\\
%\cline{2-13}
%	&F	   &	T	&	  T &	T	&	T		&	T	&	F	&	T	&	T	&	T	&	T	&	T	\\
%	&T	   &	T	&	  F &	T	&	F		&	F	&	T	&	T	&	T	&	T	&	T	&	F	\\	
%	&T	   &	T	&	 F  &	F	&	F		&	T	&	T	&	T	&	F	&	T	&	T	\\
%	&T	   &	T	&	 F  &	F	&	F		&	F	&	T	&	T	&	T	&	F	&	F	&	F	\\
%	&T	   &	F	&	F   &	T	&	T		&	T	&	T	&	F	&	T	&	T	&	T	&	T	\\
%	&T	   &	F	&	F   &	T	&	F		&	F	&	T	&	F	&	T	&	T	&	T	&	F	\\
%	&T	   &	F	&	F   &	F	&	F		&	T	&	T	&	F	&	T	&	F	&	T	&	T	\\
%	&T	   &	F	&	F   &	F	&	F		&	F	&	F	&	F	&	F	&	F	&	F	&	F	\\
%\cline{8-8}
%\end{tabular}

%\vspace{1em}

\end{compactlist}

\problempart
Escreva tabelas-verdade completas para as sentenças seguintes e marque a coluna que representa os possíveis valores de verdade da sentença inteira.

\begin{compactlist}

\item	$(D \eand \enot D) \eif G $

%\vspace{1em}

%\begin{tabular}{ccccc|c|c}
%\cline{6-6}
%1.	&	(D 	&	 \eand 	& 	 \enot	&	 D) 	&	 \eif 	&	 G \\
%	&	T	&	F		&	F		&	T	&	T	&	T	\\
%	&	T	&	F		&	F		&	T	&	T	&	F	\\
%	&	F		&	F		&	T	&	F	&	T	&	T	\\
%	&	F		&	F		&	T	&	F	&	T	&	F	\\
%\cline{6-6}
%\end{tabular}
%\vspace{1em}


\item	$(\enot P \eor \enot M) \eiff M $

%\begin{tabular}{cccccc|c|c}
%\cline{7-7}
%2.	&	(\enot 	&	P 	&	\eor 	&	\enot 	& 	 M) 	& 	\eiff 	&	 M \\
%	&	F		&	T	&	F	&	F		&	T	&	T	&	T	\\
%	&	F		&	T	&	T	&	T		&	F	&	F	&	F	\\
%	&	T		&	F	&	T	&	F		&	T	&	T	&	T	\\
%	&	T		&	F	&	T	&	T		&	F	&	T	&	F	\\
%\cline{7-7}
%\end{tabular}
%\vspace{1em}



\item	$\enot \enot (\enot A \eand \enot B)  $

%\begin{tabular}{c|c|cccccc}
%\cline{2-2}
%3.	&	\enot		&	 \enot 	&	(\enot 	& 	 A 	& \eand 	& 	\enot 	&	 B)  \\
%	&	F		&	T		&	F		&	T	&	F	&	F		&	T	\\
%	&	F		&	T		&	F		&	T	&	F	&	T		&	F	\\
%	&	F		&	T		&	T		&	F	&	F	&	F		&	T	\\
%	&	T		&	F		&	T		&	F	&	T	&	T		&	F	\\
%\cline{2-2}
%\end{tabular}
%\vspace{1em}



\item 	$[(D \eand R) \eif I] \eif \enot(D \eor R) $

%\begin{tabular}{cccccc|c|cccc}
%\cline{7-7}
%4.	&	[(D 	& 	 \eand 	& 	 R)	& 	\eif 	&	I] 	&	\eif 	&	 \enot 	&	(D 	&	 \eor 	& R) \\
%	&	T	&	T		&	T	&	T	&	T	&	F	&	F		&	T	&	T		&T	\\
%	&	T	&	T		&	T	&	F	&	F	&	T	&	F		&	T	&	T		&T	\\
%	&	T	&	F		&	F	&	T	&	T	&	F	&	F		&	T	&	T		&F	\\
%	&	T	&	F		&	F	&	T	&	F	&	F	&	F		&	T	&	T		&F	\\
%	&	F		&	F		&	T	&	T	&	T	&	F	&	F		&	F	&	T		&T	\\
%	&	F		&	F		&	T	&	T	&	F	&	F	&	F		&	F	&	T		&T	\\
%	&	F		&	F		&	F	&	T	&	T	&	T	&	T		&	F	&	F		&F	\\
%	&	F		&	F		&	F	&	T	&	F	&	T	&	T		&	F	&	F		&F	\\
%\cline{7-7}
%\end{tabular}
%	
%\vspace{1em}


\item	$\enot [(D \eiff O) \eiff A] \eif (\enot D \eand O) $

%\begin{tabular}{ccccccc|c|cccc}
%\cline{8-8}
%5.	&	\enot 	&	[(D 	&	\eiff 	&	O) 	&	\eiff 	&	 A]	& 	\eif 	 &	(\enot 	& 	D 	 & 	 \eand &O) \\ 
%	&	F		&	T	&	T	&	T	&	T	&	T	&	T	&	F		&	T	&	F	&T	\\
%	&	T		&	T	&	T	&	T	&	F	&	F	&	F	&	F		&	T	&	F	&T	\\
%	&	T		&	T	&	F	&	F	&	F	&	T	&	F	&	F		&	T	&	F	&F	\\
%	&	F		&	T	&	F	&	F	&	T	&	F	&	T	&	F		&	T	&	F	&F	\\
%	&	T		&	F	&	F	&	T	&	F	&	T	&	T	&	T		&	F	&	T	&T	\\
%	&	F		&	F	&	F	&	T	&	T	&	F	&	T	&	T		&	F	&	T	&T	\\
%	&	F		&	F	&	T	&	F	&	T	&	T	&	T	&	T		&	F	&	F	&F	\\
%	&	T		&	F	&	T	&	F	&	F	&	F	&	F	&	T		&	F	&	F	&F	\\
%\cline{8-8}
%\end{tabular}
%\vspace{1em}
\end{compactlist}


Se quiser prática adicional, você pode construir tabelas-verdade para qualquer uma das sentenças e argumentos dos exercícios do capítulo anterior.
\chapter{Conceitos semânticos}
\label{s:SemanticConcepts}

No capítulo anterior, introduzimos a ideia de uma valoração e mostramos como determinar o valor de verdade de qualquer sentença da LVF, em qualquer valoração, usando uma tabela-verdade. Neste capítulo, introduziremos algumas ideias relacionadas e mostraremos como usar tabelas-verdade para testar se elas se aplicam ou não.


\section{Tautologias e contradições}
Em \cref{s:BasicNotions}, explicamos \emph{verdade necessária} e \emph{falsidade necessária}. Ambas as noções têm substitutos na LVF. Começaremos com um substituto para a verdade necessária.
	\factoidbox{
		$\metav{A}$ é uma \define{tautologia} \ifeff{} se é verdadeira em toda valoração.
	}

\newglossaryentry{tautology}
{
name=tautologia,
description={Uma sentença que tem somente T na coluna sob o operador lógico principal de sua \gls{complete truth table}; uma sentença que é verdadeira em toda \gls{valuation}}
}

Podemos usar tabelas-verdade para decidir se uma sentença é uma tautologia. Se a sentença é verdadeira em todas as linhas de sua tabela-verdade completa, então é verdadeira em toda valoração e, portanto, é uma tautologia. No exemplo de \cref{s:CompleteTruthTables}, \textquotedblleft{}$(H \eand I) \eif H$\textquotedblright{} é uma tautologia.

No entanto, este é apenas um \emph{substituto} para a verdade necessária. Há algumas verdades necessárias que não podemos simbolizar adequadamente na LVF. Um exemplo é \textquotedblleft{}$2 + 2 = 4$\textquotedblright{}. Isso \emph{deve} ser verdadeiro, mas, se tentarmos simbolizá-lo na LVF, o melhor que podemos oferecer é uma letra sentencial, e nenhuma letra sentencial é uma tautologia. Ainda assim, se pudermos simbolizar adequadamente alguma sentença em português usando uma sentença da LVF que seja uma tautologia, então essa sentença em português expressa uma verdade necessária.

Temos um substituto semelhante para a falsidade necessária:
	\factoidbox{
		$\metav{A}$ é uma \define{contradi\c{c}\~{a}o} (na LVF) \ifeff{} é falsa em toda valoração.
	}
\newglossaryentry{contradiction of TFL}
{
  name=contradição (da LVF),
  text = contradição,
description={Uma sentença que tem somente F na coluna sob o operador lógico principal de sua \gls{complete truth table}; uma sentença que é falsa em toda \gls{valuation}}
}

Podemos usar tabelas-verdade para decidir se uma sentença é uma contradição. Se a sentença é falsa em todas as linhas de sua tabela-verdade completa, então é falsa em toda valoração e, portanto, é uma contradição. No exemplo de \cref{s:CompleteTruthTables}, \textquotedblleft{}$[(C\eiff C) \eif C] \eand \enot(C \eif C)$\textquotedblright{} é uma contradição.

\section{Equivalência}\label{sec:equivalent}
Eis uma noção útil semelhante:
	\factoidbox{
		$\metav{A}$ e $\metav{B}$ são \define{equivalentes} (na LVF) se e somente se, para toda valoração, seus valores de verdade coincidem, isto é, se não há valoração na qual tenham valores de verdade opostos.
	}
\newglossaryentry{equivalent}
{
  name=equivalência (na LVF),
  text = equivalente,
description={Uma propriedade de pares de sentenças se, e somente se, a \gls{complete truth table} dessas sentenças tem colunas idênticas sob os dois operadores lógicos principais, isto é, se as sentenças têm o mesmo valor de verdade em toda valoração}
}

Já usamos essa noção, na prática, em \cref{s:MoreBracketingConventions}; o ponto era que \textquotedblleft{}$(A \eand B) \eand C$\textquotedblright{} e \textquotedblleft{}$A \eand (B \eand C)$\textquotedblright{} são equivalentes. Novamente, é fácil testar a equivalência usando tabelas-verdade. Considere as sentenças \textquotedblleft{}$\enot(P \eor Q)$\textquotedblright{} e \textquotedblleft{}$\enot P \eand \enot Q$\textquotedblright{}. Elas são equivalentes? Para descobrir, construímos uma tabela-verdade.
\begin{center}
\begin{tabular}{c c|d e e f |d e e e f}
$P$&$Q$&\enot&$(P$&\eor&$Q)$&\enot&$P$&\eand&\enot&$Q$\\
\hline
 T & T & \TTbf{F} & T & T & T & F & T & \TTbf{F} & F & T\\
 T & F & \TTbf{F} & T & T & F & F & T & \TTbf{F} & T & F\\
 F & T & \TTbf{F} & F & T & T & T & F & \TTbf{F} & F & T\\
 F & F & \TTbf{T} & F & F & F & T & F & \TTbf{T} & T & F
\end{tabular}
\end{center}
Observe as colunas dos operadores lógicos principais: negação para a primeira sentença, conjunção para a segunda. Percorra as linhas da tabela uma a uma e compare os valores de verdade nas colunas dos operadores lógicos principais. Nas três primeiras linhas, ambos são falsos. Na última linha, ambos são verdadeiros. Como coincidem em todas as linhas, as duas sentenças são equivalentes.


\section{Satisfatibilidade}
Em \cref{s:BasicNotions}, dissemos que sentenças são conjuntamente possíveis \ifeff{} é possível que todas sejam verdadeiras ao mesmo tempo. Podemos oferecer um substituto também para essa noção:
	\factoidbox{
		$\metav{A}_1, \metav{A}_2, \ldots, \metav{A}_n$ são \define{conjuntamente satisfat\'{i}veis} (na LVF) \ifeff{} há alguma valoração que as torna todas verdadeiras.
	}

        \newglossaryentry{satisfiability in TFL}
{
  name=satisfatibilidade (na LVF),
  text=conjuntamente satisfatíveis,
  description={Uma propriedade de \glspl{sentence of TFL} se e somente se a
  \gls{complete truth table} de todas essas sentenças conjuntamente contém uma linha
  na qual todas as sentenças são verdadeiras, isto é, se alguma \gls{valuation}
  torna todas as sentenças verdadeiras}
}

Como consequência, sentenças são \define{conjuntamente insatisfat\'{i}veis} \ifeff{} nenhuma valoração torna todas verdadeiras. Novamente, é fácil testar a satisfatibilidade conjunta usando tabelas-verdade.

\section{Implicação e validade}
A ideia seguinte está intimamente relacionada à da satisfatibilidade conjunta:
	\factoidbox{
		As sentenças $\metav{A}_1, \metav{A}_2, \ldots, \metav{A}_n$ \define{implicam} (na LVF) a sentença $\metav{C}$ \ifeff{} nenhuma valoração das letras sentenciais relevantes torna todas as $\metav{A}_1, \metav{A}_2, \ldots, \metav{A}_n$ verdadeiras e $\metav{C}$ falsa.
	}
       \newglossaryentry{valid in TFL}
{
  name= validade de argumentos (na LVF),
  text = válido,
description={Uma propriedade de argumentos se, e somente se, a \gls{complete truth table} para o argumento não contém linhas nas quais as \glspl{premise} sejam todas verdadeiras e a \gls{conclusion} seja falsa, isto é, se nenhuma \gls{valuation} torna todas as premissas verdadeiras e a conclusão falsa}
}
 
Novamente, é fácil testar isso com uma tabela-verdade. Para verificar se \textquotedblleft{}$\enot L \eif (J \eor L)$\textquotedblright{} e \textquotedblleft{}$\enot L$\textquotedblright{} implicam \textquotedblleft{}$J$\textquotedblright{}, basta verificar se há alguma valoração que torne verdadeiras tanto \textquotedblleft{}$\enot L \eif (J \eor L)$\textquotedblright{} quanto \textquotedblleft{}$\enot L$\textquotedblright{}, ao mesmo tempo que torna \textquotedblleft{}$J$\textquotedblright{} falsa. Portanto, usamos uma tabela-verdade:
\begin{center}
\begin{tabular}{c c|d e e e e f|d f| c}
$J$&$L$&\enot&$L$&\eif&$(J$&\eor&$L)$&\enot&$L$&$J$\\
\hline
%J   L   -   L      ->     (J   v   L)
 T & T & F & T & \TTbf{T} & T & T & T & \TTbf{F} & T & \TTbf{T}\\
 T & F & T & F & \TTbf{T} & T & T & F & \TTbf{T} & F & \TTbf{T}\\
 F & T & F & T & \TTbf{T} & F & T & T & \TTbf{F} & T & \TTbf{F}\\
 F & F & T & F & \TTbf{F} & F & F & F & \TTbf{T} & F & \TTbf{F}
\end{tabular}
\end{center}
A única linha na qual tanto \textquotedblleft{}$\enot L \eif (J \eor L)$\textquotedblright{} quanto \textquotedblleft{}$\enot L$\textquotedblright{} são verdadeiras é a segunda, e essa é uma linha na qual \textquotedblleft{}$J$\textquotedblright{} também é verdadeira. Portanto, \textquotedblleft{}$\enot L \eif (J \eor L)$\textquotedblright{} e \textquotedblleft{}$\enot L$\textquotedblright{} implicam \textquotedblleft{}$J$\textquotedblright{}.

% We now make an important observation:
% 	\factoidbox{
% 	If $\metav{A}_1, \metav{A}_2, \ldots, \metav{A}_n$ entail
% 	$\metav{C}$ in TFL, then $\metav{A}_1, \metav{A}_2, \ldots,
% 	\metav{A}_n \therefore \metav{C}$ is valid.
% 	}
% Here's why. Suppose $\metav{A}_1, \metav{A}_2, \ldots, \metav{A}_n$
% entail $\metav{C}$ in TFL. Can it happen that the argument
% $\metav{A}_1, \metav{A}_2, \ldots, \metav{A}_n \therefore \metav{C}$
% is invalid? Then there would be a \emph{case} which makes all of
% $\metav{A}_1, \metav{A}_2, \ldots, \metav{A}_n$ true and $\metav{C}$
% false, relative to some symbolization key. This case would generate a
% valuation of the sentence letters occurring in $\metav{A}_1,
% \metav{A}_2, \ldots, \metav{A}_n$, and $\metav{C}$: take the truth
% value of any sentence letter to be just the truth value of the
% corresponding sentence in the case in question. This valuation would
% also make $\metav{A}_1, \metav{A}_2, \ldots, \metav{A}_n$ true and
% $\metav{C}$ false, since truth values are determined from the truth
% values of sentence letters by the truth tables of the connectives, and
% the truth values of the sentence letters are the same in the valuation
% as they are in the (imagined) case. But this is impossible, since
% we've assumed that $\metav{A}_1, \metav{A}_2, \ldots, \metav{A}_n$
% entail $\metav{C}$ in TFL, and so there is no valuation which makes
% all of $\metav{A}_1, \metav{A}_2, \ldots, \metav{A}_n$ true and also
% makes $\metav{C}$ false. So it can't happen that the argument
% $\metav{A}_1, \metav{A}_2, \ldots, \metav{A}_n \therefore \metav{C}$
% is invalid; consequently, it must be valid.

% In short, we have a way to test for the validity of English arguments. First, we symbolize them in TFL; then we test for entailment in TFL using truth tables.



\section{O duplo turnstile}
Neste livro, usaremos bastante a noção de implicação. Será útil, então, introduzir um símbolo que a abrevia. Em vez de dizer que as sentenças da LVF $\metav{A}_1$, $\metav{A}_2$, \dots{} e $\metav{A}_n$ conjuntamente implicam $\metav{C}$, abreviaremos isso assim:
	$$\metav{A}_1, \metav{A}_2, \dots, \metav{A}_n \entails \metav{C}$$
O símbolo \textquotedblleft{}$\entails$\textquotedblright{} é conhecido como \emph{duplo turnstile}, pois se parece com um turnstile com duas barras horizontais.

Sejamos claros. \textquotedblleft{}$\entails$\textquotedblright{} não é um símbolo da LVF. Antes, é um símbolo de nossa metalinguagem, o português ampliado (recorde a diferença entre linguagem-objeto e metalinguagem em \cref{s:UseMention}). Assim, a sentença da metalinguagem:
	\[\metav{A}, \metav{A} \eif \metav{B} \entails \metav{B}\]
é \emph{apenas} uma abreviação desta sentença da metalinguagem:
	\begin{center}
		As sentenças da LVF $\metav{A}$ e $\metav{A} \eif \metav{B}$ implicam $\metav{B}$
	\end{center}
Observe que não há limite para o número de sentenças da LVF que podem ser mencionadas antes do símbolo~\textquotedblleft{}$\entails$\textquotedblright{}. Com efeito, podemos até considerar o caso-limite:
	$$\entails \metav{C}$$
Isso diz que não há valoração que torne verdadeiras todas as sentenças mencionadas no lado esquerdo de \textquotedblleft{}$\entails$\textquotedblright{} e, ao mesmo tempo, torne falsas todas as sentenças no lado direito (neste caso, $\metav{C}$). Como \emph{nenhuma} sentença é mencionada no lado esquerdo de \textquotedblleft{}$\entails$\textquotedblright{} neste caso, isso significa simplesmente que não há valoração que torne $\metav{C}$ falsa. Dito de outro modo, significa que toda valoração torna $\metav{C}$ verdadeira. Dito de outro modo, significa que $\metav{C}$ é uma tautologia. Da mesma forma, para dizer que $\metav{A}$ é uma contradição, poderíamos escrever:
	$$\metav{A} \entails\phantom{\metav{C}}$$
para significar que nenhuma valoração torna $\metav{A}$ verdadeira.

Às vezes, queremos negar que haja uma implicação tautológica e dizer algo desta forma:
\begin{center}
	não é o caso que $\metav{A}_1, \ldots, \metav{A}_n \entails \metav{C}$
\end{center}
Nesse caso, podemos simplesmente riscar o turnstile e escrever:
$$\metav{A}_1, \metav{A}_2, \ldots, \metav{A}_n \nentails\metav{C}$$
Isso significa que alguma valoração torna todas as $\metav{A}_1, \ldots, \metav{A}_n$ verdadeiras e torna $\metav{C}$ falsa. Observe que disso \emph{não} se segue que $\metav{A}_1,\ldots, \metav{A}_n \entails \enot \metav{C}$, pois é possível que alguma \emph{outra} valoração torne todas as $\metav{A}_1, \ldots, \metav{A}_n$ verdadeiras e torne $\metav{C}$ \emph{verdadeira}. Por exemplo, $P \nentails Q$, mas também $P \nentails \enot Q$.

\section{$\entails$ versus $\eif$}\label{entails-v-iff}
Agora queremos comparar e contrastar \textquotedblleft{}$\entails$\textquotedblright{} e \textquotedblleft{}$\eif$\textquotedblright{}.

Observe: $\metav{A} \entails \metav{C}$ \ifeff{} nenhuma valoração das letras sentenciais torna $\metav{A}$ verdadeira e $\metav{C}$ falsa.

Observe: $\metav{A} \eif \metav{C}$ é uma tautologia \ifeff{} nenhuma valoração das letras sentenciais torna $\metav{A} \eif \metav{C}$ falsa. Como um condicional é verdadeiro exceto quando seu antecedente é verdadeiro e seu consequente é falso, $\metav{A} \eif \metav{C}$ é uma tautologia \ifeff{} nenhuma valoração torna $\metav{A}$ verdadeira e $\metav{C}$ falsa.

Combinando essas duas observações, vemos que $\metav{A} \eif \metav{C}$ é uma tautologia \ifeff{} $\metav{A} \entails \metav{C}$. Mas há uma diferença realmente, realmente importante entre \textquotedblleft{}$\entails$\textquotedblright{} e \textquotedblleft{}$\eif$\textquotedblright{}:
	\factoidbox{
		\begin{itemize}
			\item \textquotedblleft{}$\eif$\textquotedblright{} é um conectivo sentencial da LVF.
			\item \textquotedblleft{}$\entails$\textquotedblright{} é um símbolo do português ampliado.
		\end{itemize}
	}
Com efeito, quando \textquotedblleft{}$\eif$\textquotedblright{} é ladeado por duas sentenças da LVF, o resultado é uma sentença mais longa da LVF. Em contraste, quando usamos \textquotedblleft{}$\entails$\textquotedblright{}, formamos uma sentença metalinguística que \emph{menciona} as sentenças da LVF ao seu redor.

A palavra inglesa \textquotedblleft implies\textquotedblright{} é frequentemente usada como sinônimo de \textquotedblleft entails\textquotedblright{}, isto é, de \textquotedblleft{}$\entails$\textquotedblright{}. Alguns lógicos \emph{também} a usam para o condicional, isto é, para \textquotedblleft{}$\eif$\textquotedblright{}. Como isso pode levar à confusão alertada acima, evitamos usar \textquotedblleft implies\textquotedblright{} tanto quanto possível. (Quando ela é usada, significa \textquotedblleft entails\textquotedblright{}.)

\practiceproblems
\problempart
Reveja suas respostas para
\hyperref[pr.TT.TTorC]{exercício~\ref*{s:CompleteTruthTables}A}. Determine quais
sentenças eram tautologias, quais eram contradições e quais eram
não eram nem tautologias nem contradições.
\solutions

\

\problempart
\label{pr.TT.satisfiable}
Use tabelas-verdade para determinar se estas sentenças são conjuntamente satisfatíveis ou conjuntamente insatisfatíveis:
\begin{compactlist}
\item $A\eif A$, $\enot A \eif \enot A$, $A\eand A$, $A\eor A$ %satisfiable
\item $A\eor B$, $A\eif C$, $B\eif C$ %satisfiable
\item $B\eand(C\eor A)$, $A\eif B$, $\enot(B\eor C)$  %unsatisfiable
\item $A\eiff(B\eor C)$, $C\eif \enot A$, $A\eif \enot B$ %satisfiable
\end{compactlist}


\solutions
\problempart
\label{pr.TT.valid}
Use tabelas-verdade para determinar se cada argumento é válido ou inválido.
\begin{compactlist}
\item $A\eif A \therefore A$ %invalid
\item $A\eif(A\eand\enot A) \therefore \enot A$ %valid
\item $A\eor(B\eif A) \therefore \enot A \eif \enot B$ %valid
\item $A\eor B, B\eor C, \enot A \therefore B \eand C$ %invalid
\item $(B\eand A)\eif C, (C\eand A)\eif B \therefore (C\eand B)\eif A$ %invalid
\end{compactlist}

\problempart Determine se cada sentença é uma tautologia, uma contradição ou uma sentença contingente, usando uma tabela-verdade completa.
\begin{compactlist}
\item $\enot B \eand B$ \vspace{.5ex}%contra


\item $\enot D \eor D$ \vspace{.5ex}%taut


\item $(A\eand B) \eor (B\eand A)$\vspace{.5ex} %contingent


\item $\enot[A \eif (B \eif A)]$\vspace{.5ex} %contra


\item $A \eiff [A \eif (B \eand \enot B)]$ \vspace{.5ex}%contra


\item $[(A \eand B) \eiff B] \eif (A \eif B)$ \vspace{.5ex}% contingent.

\end{compactlist}



\noindent\problempart
\label{pr.TT.equiv}
Determine se as sentenças seguintes são logicamente equivalentes usando tabelas-verdade completas. Se as duas sentenças forem realmente logicamente equivalentes, escreva ``equivalentes.'' Caso contrário, escreva ``não equivalentes.'' 
\begin{compactlist}
\item $A$ e $\enot A$
\item $A \eand \enot A$ e $\enot B \eiff B$
\item $[(A \eor B) \eor C]$ e $[A \eor (B \eor C)]$
\item $A \eor (B \eand C)$ e $(A \eor B) \eand (A \eor C)$
\item $[A \eand (A \eor B)] \eif B$ e $A \eif B$\end{compactlist}


\problempart
\label{pr.TT.equiv2}
Determine se as sentenças seguintes são logicamente equivalentes usando tabelas-verdade completas. Se as duas sentenças forem realmente equivalentes, escreva ``equivalentes.'' Caso contrário, escreva ``não equivalentes.''
\begin{compactlist}
\item $A\eif A$ e $A \eiff A$
\item $\enot(A \eif B)$ e $\enot A \eif \enot B$
\item $A \eor B$ e $\enot A \eif B$
\item$(A \eif B) \eif C$ e $A \eif (B \eif C)$
\item $A \eiff (B \eiff C)$ e $A \eand (B \eand C)$
\end{compactlist}


\problempart
\label{pr.TT.satisfiable2}
Determine se cada coleção de sentenças é conjuntamente satisfatível ou conjuntamente insatisfatível usando uma tabela-verdade completa.
\begin{compactlist}
\item $A \eand \enot B$, $\enot(A \eif B)$, $B \eif A$\vspace{.5ex} %Consistent

%\begin{tabular}{ccccccccccccccc} 
%1. 	&	A 					 & \eand 		&  \enot & B & & \enot  		& 	 (A	  & 	 \eif	 	 & 	 B)		 & 	 & 	 B	 	 & 	\eif 	 	 & 	A 	 	 & 	 Consistent \\ 
%\cline{2-5} \cline{7-10}\cline{12-14} 
%	& 	T 					 & 	 F	 		&  F	 & T & & F	 		& 	 T	  & 	 T	 	 & 	T 	 	 & 	 & 	 T	 	 & 	 T	 	 & T	 	 	&	  \\ 
%\cline{2-14}
%	& \multicolumn{1}{|r}{T}& 	\textbf{T}	 & T	 & F & & \textbf{T}	 & 	 T	 & 	 F	 	 & 	 F	 	 & 	 & 	 F	 	 & 	 \textbf{T}	 	 & 	 \multicolumn{1}{r|}{T}	 	 & 	  \\ 
%\cline{2-14}
%	& 	 F	 				 & 	 F	 & 	 F	 & T & 	& 	 F	 & 	 F	 & 	 T	 	 & 	 T	 	 & 	  & 	 T	 	 & 	 F	 	 & 	 F	 	 & 	  \\ 
%	& 	 F	  				& 	 F	 & 	 T	 & 	F&  & 	 F	 & 	 F	 & 	 T	 	 & 	 F	 	 & 	  & 	 F	 	 & 	 T	 	 & 	 F	 	 & 	  \\ 
%\end{tabular}

\item $A \eor B$, $A \eif \enot A$, $B \eif \enot B$ \vspace{.5ex}%unsatisfiable.

%\begin{tabular}{ccccccccccccccc} 
%2. &A	 & \eor 	 & B 	 & 	 	 & A 	 & \eif 	 & 	\enot & A 	 & 	 	 & B 	 & \eif 	 & \enot	 & 	B 	 & 	Insatisfiable \\ 
%\cline{2-4}\cline{6- 9} \cline{11-14}
%   &	T	 & 	 T	 &T  	 & 	 	 & T	 & 	 F	 & 	F 	 & T 	 & 	 	 & 	T 	 & 	F 	 & 	 F	 & 	T 	 & 	 \\ 
%   &	 T	& 	 T	 & F 	 & 	 	 & 	T 	 & 	 F	 & 	 F	 & 	 T	 & 	 	 & 	F 	 & 	 T	 & 	 T	 & 	 F	 & 	 \\ 
%   &	 F	& 	 T	 & 	 T	 & 	 	 & 	F 	 & 	 T	 & 	 T	 & 	F 	 & 	 	 & 	 T	 & 	 F	 & 	 F	 & 	 T	 & 	 \\ 
%   &	 F	& 	 F	 & 	 F	 & 	 	 & 	 F	 & 	 T	 & 	 T	 & 	 F	 & 	 	 & 	 F	 & 	 T	 & 	 T	 & 	 F	 & 	 \\ 
%\end{tabular}

\item $\enot(\enot A \eor B) $, $A \eif \enot C$, $A \eif (B \eif C)$\vspace{.5ex} %Insatisfiable

%3. &\enot & (\enot & A & \eor &B) &  &A  & \eif 	 &\enot 	 &C & 	 & A &\eif 	& (B 	 &\eif 	& C)	 &Consistent \\ 
%\cline{2-6}\cline{8-11} \cline{13-17} 
%   &	F 	& 	F	 & 	T & T	 & T & 	  & T & F	 & 	 F&T 	 & 	 &T & T	 & T	 &T 	 &T 	 & \\ 
%   &	 F	& 	F	 & 	T & T	 & T & 	  & T & T	 & 	 T& F	 & 	 &T & F	 & T	 & F	 &F 	 & \\ 
% 
%  &	 T & 	F 	& 	T & F	 & F & 	  & T & F	 & 	 F& T	 & 	 &T & T	 & F	 & T	 &T 	 & \\ 
%\cline{2-17}
%   &	 \multicolumn{1}{|r}{{\color{red}T}}		&  F	 & 	T & F	 & 	F &  & 	T & {\color{red}T}	 & 	 T&F 	& 	 &T & {\color{red}T}	 & F	 & T	 &\multicolumn{1}{r|}{F} 	 & \\ 
%\cline{2-17}
%   &	 F	& 	T	 & 	F & T	 & 	T &  & 	F & T	 & 	 F& T	 & 	 &F	 & F	 & T	 & T	 &T 	 & \\ 
%   &	 F	& 	 T	& 	F & T	 & 	T &  & 	F & T	 & 	T & F 	& 	 &F	 & T	 & T	 &F 	 &F 	 & \\ 
%   &	 F	& 	 T	& 	F & T	 & 	F &  & 	F & T	 & 	F & T	 & 	 &F	 & T	 & F	 & T	 &T 	 & \\ 
%   &	 F	& 	 T	& 	F & T	 & 	F &  & 	F & T	 & 	T & F	 & 	 &F	 & T	 & F	 & T	 &F 	 & \\ 
%\end{tabular}
%


\item $A \eif B$, $A \eand \enot B$\vspace{.5ex} %Insatisfiable

\item $A \eif (B \eif C)$, $(A \eif B) \eif C$, $A \eif C$\vspace{.5ex} % satisfiable.

\end{compactlist}

\noindent\problempart
\label{pr.TT.satisfiable3}
Determine se cada coleção de sentenças é conjuntamente satisfatível ou conjuntamente insatisfatível, usando uma tabela-verdade completa.
\begin{compactlist}
\item $\enot B$, $A \eif B$, $A$ \vspace{.5ex}%unsatisfiable.
\item $\enot(A \eor B)$, $A \eiff B$, $B \eif A$\vspace{.5ex} %Consistent
\item $A \eor B$, $\enot B$, $\enot B \eif \enot A$\vspace{.5ex} %Insatisfiable
\item $A \eiff B$, $\enot B \eor \enot A$, $A \eif B$\vspace{.5ex} %satisfiable.
\item $(A \eor B) \eor C$, $\enot A \eor \enot B$, $\enot C \eor \enot B$\vspace{.5ex} %satisfiable
\end{compactlist}




\noindent\problempart
\label{pr.TT.valid2}
Determine se cada argumento é válido ou inválido, usando uma tabela-verdade completa.
\begin{compactlist}
\item $A\eif B$, $B \therefore  A$ %invalid

\item $A\eiff B$, $B\eiff C \therefore A\eiff C$ %valid

\item $A \eif B$, $A \eif C\therefore B \eif C$ %invalid.

\item $A \eif B$, $B \eif A\therefore A \eiff B$ %valid.

\end{compactlist}

\noindent\problempart
\label{pr.TT.valid3}
Determine se cada argumento é válido ou inválido, usando uma tabela-verdade completa.
\begin{compactlist}
\item $A\eor\bigl[A\eif(A\eiff A)\bigr] \therefore  A $\vspace{.5ex}%invalid
\item $A\eor B$, $B\eor C$, $\enot B \therefore A \eand C$\vspace{.5ex} %valid
\item $A \eif B$, $\enot A\therefore \enot B$ \vspace{.5ex}%invalid
\item $A$, $B\therefore \enot(A\eif \enot B)$ \vspace{.5ex}%valid
\item $\enot(A \eand B)$, $A \eor B$, $A \eiff B\therefore C$ \vspace{.5ex}%valid 
\end{compactlist}

\solutions
\problempart
\label{pr.TT.concepts}
Responda a cada uma das perguntas abaixo e justifique sua resposta.
\begin{compactlist}
\item Suponha que \metav{A} e \metav{B} sejam logicamente equivalentes. O que se pode dizer sobre $\metav{A}\eiff\metav{B}$?
%\metav{A} and \metav{B} have the same truth value on every line of a complete truth table, so $\metav{A}\eiff\metav{B}$ is true on every line. It is a tautology.
\item Suponha que $(\metav{A}\eand\metav{B})\eif\metav{C}$ não seja nem uma tautologia nem uma contradição. O que se pode dizer sobre a validade de $\metav{A}, \metav{B} \therefore\metav{C}$?
%The sentence is false on some line of a complete truth table. On that line, \metav{A} and \metav{B} are true and \metav{C} is false. So the argument is invalid.
\item Suponha que $\metav{A}$, $\metav{B}$ e $\metav{C}$  sejam conjuntamente insatisfatíveis. O que se pode dizer sobre $(\metav{A}\eand\metav{B}\eand\metav{C})$?
\item Suponha que \metav{A} seja uma contradição. O que se pode dizer sobre a validade de $\metav{A}, \metav{B} \entails \metav{C}$?
%Since \metav{A} is false on every line of a complete truth table, there is no line on which \metav{A} and \metav{B} are true and \metav{C} is false. So the argument is valid.
\item Suponha que \metav{C} seja uma tautologia. O que se pode dizer sobre a validade de $\metav{A}, \metav{B}\entails \metav{C}$?
%Since \metav{C} is true on every line of a complete truth table, there is no line on which \metav{A} and \metav{B} are true and \metav{C} is false. So the argument is valid.
\item Suponha que \metav{A} e \metav{B} sejam logicamente equivalentes. O que se pode dizer sobre $(\metav{A}\eor\metav{B})$?
%Not much. $(\metav{A}\eor\metav{B})$ is a tautology if \metav{A} and \metav{B} are tautologies; it is a contradiction if they are contradictions; it is contingent if they are contingent.
\item Suponha que \metav{A} e \metav{B} não sejam logicamente equivalentes. O que se pode dizer sobre $(\metav{A}\eor\metav{B})$?
%\metav{A} and \metav{B} have different truth values on at least one line of a complete truth table, and $(\metav{A}\eor\metav{B})$ will be true on that line. On other lines, it might be true or false. So $(\metav{A}\eor\metav{B})$ is either a tautology or it is contingent; it is \emph{not} a contradiction.
\end{compactlist}
\problempart 
Considere o seguinte princípio:
	\begin{quote}
		Suponha que $\metav{A}$ e $\metav{B}$ sejam logicamente equivalentes. Suponha que um argumento contenha $\metav{A}$ (seja como premissa, seja como conclusão). A validade do argumento não seria afetada se substituíssemos $\metav{A}$ por $\metav{B}$.
	\end{quote}
Esse princípio está correto? Explique sua resposta.

\chapter{Limitações da LVF}\label{s:ParadoxesOfMaterialConditional}

Alcançamos um marco importante: um teste para a validade de
argumentos! No entanto, ainda não devemos nos empolgar. É importante
compreender os \emph{limites} da nossa conquista. Ilustraremos esses
limites com quatro exemplos.

Primeiro, considere o argumento:
	\begin{earg}
		\item Daisy tem quatro patas.
		\item[\texttherefore] Daisy tem mais de duas patas.
	\end{earg}
Para simbolizar esse argumento na LVF, teríamos de usar duas letras
sentenciais diferentes---talvez `$F$' e `$T$'---para a premissa e a
conclusão, respectivamente. Ora, é evidente que `$F$' não implica~`$T$'.
Mas o argumento em português é certamente válido!

Segundo, considere a sentença:
	\begin{enumerate}
		\item\label{n:JanBald} Jan não é nem careca nem não careca.
	\end{enumerate}
Para simbolizar essa sentença na LVF, ofereceríamos algo como
`$\enot J \eand \enot \enot J$'. Isso é uma contradição (verifique com
uma tabela-verdade), mas a sentença \cref*{n:JanBald} em si não parece
ser uma contradição: poderíamos perfeitamente acrescentar `Jan está no
limiar da calvície'!

Terceiro, considere a seguinte sentença:
	\begin{enumerate}
	\item\label{n:GodParadox} Não é o caso de que, se Deus existe, Ela atende a preces malévolas.
	\end{enumerate}
Ao simbolizá-la na LVF, ofereceríamos algo como `$\enot (G \eif M)$'.
Ora, `$\enot (G \eif M)$' implica `$G$' (novamente, verifique com uma
tabela-verdade). Assim, se simbolizarmos \cref*{n:GodParadox} na LVF,
parece que ela implica que Deus existe. Mas isso é estranho: certamente
até mesmo um ateu pode aceitar \cref*{n:GodParadox} sem se contradizer!

Uma lição disso é que a simbolização de \cref*{n:GodParadox} como
`$\enot(G \eif M)$' não expressa o que pretendemos. Talvez devêssemos
reformular \cref*{n:GodParadox} assim:
	\begin{enumerate}
	\item\label{n:GodParadox2} Se Deus existe, Ela não atende a preces malévolas.
	\end{enumerate}
e simbolizar \cref*{n:GodParadox2} como `$G \eif \enot M$'. Agora, se os
ateus estiverem certos e não houver Deus, então `$G$' é falsa e, portanto,
`$G \eif \enot M$' é verdadeira, e o problema desaparece. No entanto,
se `$G$' é falsa, `$G \eif M$', isto é, `Se Deus existe, Ela atende a
preces malévolas', também é \emph{verdadeira}!

De maneiras diferentes, esses quatro exemplos destacam alguns dos
limites da LVF, da simbolização do português na LVF e dos testes baseados
em tabelas-verdade que elaboramos. Nosso primeiro exemplo mostrou que a
LVF não é suficientemente expressiva para simbolizar tudo o que podemos
querer, de um modo que nos permita aplicar nosso instrumental lógico até
mesmo à melhor simbolização possível. Veremos mais adiante
(\cref{sec.identity}) que \emph{podemos} simbolizar corretamente `Daisy
tem quatro patas' na linguagem mais expressiva LPO, e então o teste que
elaboraremos para a LPO será aplicável (e dará a resposta correta).

Em \cref{n:JanBald}, tivemos de modo semelhante um exemplo em que a
simbolização direta na LVF não funciona inteiramente. O melhor que podemos
fazer na LVF é simbolizar `Jan não é careca' usando sua própria letra
sentencial, digamos,~`$N$' (contrariamente à nossa prática usual). Mas isso
também não funciona. Por exemplo, então simbolizaríamos `Jan é careca e
não é' como `$J \eand N$'. Mas `Jan é careca e não é' é, plausivelmente,
uma autocontradição, enquanto a simbolização alternativa `$J \eand N$' não
é. Alguns lógicos propuseram que, nos casos de sentenças como `Jan é
careca', que admitem casos limítrofes, devemos ajustar nossa semântica e
permitir que a letra sentencial `$J$' que a simboliza possa assumir valores
de verdade diferentes de T e~F. Mas isso está longe de ser universalmente
aceito.

\Cref{n:GodParadox} também mostrou que a simbolização costuma ser
complicada e que a simbolização direta de uma sentença em português às
vezes não capta suas condições de verdade pretendidas. Mas até mesmo nossa
melhoria, \cref{n:GodParadox2}, e sua simbolização `$G \eif
\enot M$' acabaram não sendo suficientes. O fenômeno que encontramos aqui
é um dos chamados paradoxos do condicional material. O paradoxo é que
condicionais materiais como `$G \eif M$' e `$G \eif \enot M$' são
verdadeiros sempre que seu antecedente~`$G$' é falso, embora, intuitivamente,
as duas sentenças `Se Deus existe, Ela atende a preces malévolas' e `Se Deus
existe, Ela \emph{não} atende a preces malévolas' não devessem ser ambas
verdadeiras. Isso indica que o `se \dots então' nesses dois exemplos não é
captado adequadamente pelo conectivo da LVF~`\eif'. A verofuncionalidade da
LVF é uma limitação real nesse caso. A solução seria tratar o `se \dots
então' nesse caso como um condicional subjuntivo (ver
\cref{s:IndicativeSubjunctive}), mas a LVF não consegue captar esse tipo de
condicional.\footnote{Os lógicos desenvolveram lógicas que lidam melhor com
condicionais subjuntivos. Elas são chamadas de `lógicas condicionais' ou
`lógicas dos contrafactuais'.}

O caso da calvície de Jan (ou da ausência dela) levanta a questão geral de
qual lógica devemos usar ao lidar com discurso \emph{vago}. O caso do ateu
levanta a questão de como lidar com os (assim chamados) \emph{paradoxos do
condicional material}. Um dos objetivos deste livro é fornecer as
ferramentas para explorar essas questões de \emph{lógica filosófica}. Mas
precisamos aprender a andar antes de correr; temos de nos tornar proficientes
no uso da LVF antes de podermos discutir adequadamente suas limitações e
considerar alternativas.

\chapter{Atalhos de tabelas-verdade}
Com a prática, você rapidamente se tornará hábil em preencher tabelas-verdade. Neste capítulo, consideramos (e justificamos) alguns atalhos que ajudarão ao longo do caminho.

\section{Trabalhando com tabelas-verdade}
Você logo descobrirá que não precisa copiar o valor de verdade de cada letra sentencial, mas pode simplesmente consultá-las. Assim, pode agilizar o trabalho escrevendo:
\begin{center}
\begin{tabular}{c c|d e e e e f}
$P$&$Q$&$(P$&\eor&$Q)$&\eiff&\enot&$P$\\
\hline
 T & T &  & T &  & \TTbf{F} & F\\
 T & F &  & T &  & \TTbf{F} & F\\
 F & T &  & T & & \TTbf{T} & T\\
 F & F &  & F &  & \TTbf{F} & T
\end{tabular}
\end{center}
Você também sabe com certeza que uma disjunção é verdadeira sempre que um dos disjuntivos é verdadeiro. Portanto, se encontrar um disjuntivo verdadeiro, não há necessidade de calcular o valor de verdade do outro disjuntivo. Assim, você poderia apresentar:
\begin{center}
\begin{tabular}{c c|d e e e e e e f}
$P$&$Q$& $(\enot$ & $P$&\eor&\enot&$Q)$&\eor&\enot&$P$\\
\hline
 T & T & F & & F & F& & \TTbf{F} & F\\
 T & F &  F & & T& T& &  \TTbf{T} & F\\
 F & T & & &  & & & \TTbf{T} & T\\
 F & F & & & & & &\TTbf{T} & T
\end{tabular}
\end{center}
Do mesmo modo, você sabe com certeza que uma conjunção é falsa sempre que um dos conjuntivos é falso. Portanto, se encontrar um conjuntivo falso, não há necessidade de calcular o valor de verdade do outro conjuntivo. Assim, você poderia apresentar:
\begin{center}
\begin{tabular}{c c|d e e e e e e f}
$P$&$Q$&\enot &$(P$&\eand&\enot&$Q)$&\eand&\enot&$P$\\
\hline
 T & T &  &  & &  & & \TTbf{F} & F\\
 T & F &   &  &&  & & \TTbf{F} & F\\
 F & T & T &  & F &  & & \TTbf{T} & T\\
 F & F & T &  & F & & & \TTbf{T} & T
\end{tabular}
\end{center}
Um atalho semelhante está disponível para condicionais. Você sabe imediatamente que um condicional é verdadeiro se seu consequente é verdadeiro ou se seu antecedente é falso. Assim, você poderia apresentar:
\begin{center}
\begin{tabular}{c c|d e e e e e f}
$P$&$Q$& $((P$&\eif&$Q$)&\eif&$P)$&\eif&$P$\\
\hline
 T & T & &  & & & & \TTbf{T} & \\
 T & F &  &  & && & \TTbf{T} & \\
 F & T & & T & & F & & \TTbf{T} & \\
 F & F & & T & & F & &\TTbf{T} & 
\end{tabular}
\end{center}
Portanto, `$((P \eif Q) \eif P) \eif P$' é uma tautologia. Na verdade, é uma instância da \emph{Lei de Peirce}, assim chamada em homenagem a Charles Sanders Peirce.

\section{Testando validade e implicação}
Em \cref{s:SemanticConcepts}, vimos como usar tabelas-verdade para testar a validade. Nesse teste, procuramos \emph{linhas ruins}: linhas em que todas as premissas são verdadeiras e a conclusão é falsa. Agora:
\begin{itemize}
	\item Se a conclusão é verdadeira em uma linha, então essa linha não é ruim. (E não precisamos avaliar mais nada \emph{além} disso nessa linha para confirmar o fato.)
	\item Se alguma premissa é falsa em uma linha, então essa linha não é ruim. (E não precisamos avaliar mais nada \emph{além} disso nessa linha para confirmar o fato.)
\end{itemize}
Com isso em mente, podemos acelerar consideravelmente nossos testes de validade.

Consideremos como podemos testar o seguinte:
$$\enot L \eif (J \eor L), \enot L \therefore J$$
A \emph{primeira} coisa que devemos fazer é avaliar a conclusão. Se descobrirmos que a conclusão é \emph{verdadeira} em alguma linha, então essa não é uma linha ruim. Assim, podemos simplesmente ignorar o restante da linha. Depois dessa primeira etapa, ficamos com algo como:
\begin{center}
	\begin{tabular}{c c|d e e e e f |d f|c}
		$J$&$L$&\enot&$L$&\eif&$(J$&\eor&$L)$&\enot&$L$&$J$\\
		\hline
		%J   L   -   L      ->     (J   v   L)
		T & T & &&&&&&&& {T}\\
		T & F & &&&&&&&& {T}\\
		F & T & &&?&&&&?&& {F}\\
		F & F & &&?&&&&?&& {F}
	\end{tabular}
\end{center}
em que os espaços em branco indicam que não nos daremos ao trabalho de continuar investigando (pois a linha não é ruim), e os pontos de interrogação indicam que precisamos continuar investigando.

A premissa mais fácil de avaliar é a segunda; fazemos isso em seguida e obtemos:
\begin{center}
	\begin{tabular}{c c|d e e e e f |d f|c}
		$J$&$L$&\enot&$L$&\eif&$(J$&\eor&$L)$&\enot&$L$&$J$\\
		\hline
		%J   L   -   L      ->     (J   v   L)
		T & T & &&&&&&&& {T}\\
		T & F & &&&&&&&& {T}\\
		F & T & &&&&&&{F}&& {F}\\
		F & F & &&?&&&&{T}&& {F}
	\end{tabular}
\end{center}
Observe que não precisamos mais considerar a terceira linha da tabela: ela certamente não é ruim, pois alguma premissa é falsa nessa linha. E, finalmente, completamos a tabela-verdade:
\begin{center}
	\begin{tabular}{c c|d e e e e f |d f|c}
		$J$&$L$&\enot&$L$&\eif&$(J$&\eor&$L)$&\enot&$L$&$J$\\
		\hline
		%J   L   -   L      ->     (J   v   L)
		T & T & &&&&&&&& {T}\\
		T & F & &&&&&&&& {T}\\
		F & T & &&&&&&{F}& & {F}\\
		F & F & T &  & \TTbf{F} &  & F & & {T} & & {F}
	\end{tabular}
\end{center}
A tabela-verdade não tem linhas ruins, portanto o argumento é válido. Qualquer valoração que torne verdadeiras todas as premissas torna verdadeira a conclusão.

Vale a pena ilustrar a estratégia novamente. Considere este argumento:
$$A\eor B, \enot (B\eand C) \therefore (A \eor \enot C)$$
Mais uma vez, começamos avaliando a conclusão. Como ela é uma disjunção, é verdadeira sempre que qualquer um dos disjuntivos é verdadeiro, de modo que podemos acelerar um pouco o trabalho.
\begin{center}
\begin{tabular}[t]{c c c| c|c|d e e f }
$A$ & $B$ & $C$ & $A\eor B$ & $\enot (B \eand C)$ & $(A$ &$\eor $& $\enot $ & $C)$\\
\hline
T & T & T &  &  & & \TTbf{T} & & \\
T & T & F &  &  & & \TTbf{T} & & \\
T & F & T &  &  & & \TTbf{T} & & \\
T & F & F &  &  & & \TTbf{T} & & \\
F & T & T & ? & ? & & \TTbf{F} &F & \\
F & T & F &  &  && \TTbf{T} & T& \\
F & F & T & ? & ? && \TTbf{F} & F& \\
F & F & F &  &  & & \TTbf{T} & T& \\
\end{tabular}
\end{center}
Agora podemos ignorar todas as linhas, exceto as duas em que a sentença depois do símbolo de consequência é falsa. Avaliando as duas sentenças à esquerda do símbolo de consequência, obtemos:
\begin{center}
	\begin{tabular}[t]{c c c| c|d e e f |d e e f }
		$A$ & $B$ & $C$ & $A\eor B$ & $\enot ($&$B$&$ \eand$&$ C)$ & $(A$ &$\eor $& $\enot $ & $C)$\\
		\hline
		T & T & T &  & &&& & & \TTbf{T} & & \\
		T & T & F &  & &&& & & \TTbf{T} & & \\
		T & F & T &  & &&& & & \TTbf{T} & & \\
		T & F & F &  & &&& & & \TTbf{T} & & \\
		F & T & T & \textbf{T} & \textbf{F}&&T& & & \TTbf{F} &F & \\
		F & T & F & &&& & && \TTbf{T} & T& \\
		F & F & T & \textbf{F} & &&& & & \TTbf{F} & F& \\
		F & F & F & &&&& && \TTbf{T} & T& \\
	\end{tabular}
\end{center}
Logo, a implicação vale! E nossos atalhos nos pouparam \emph{muito} trabalho.
 
Estivemos discutindo atalhos para testar a validade. Mas exatamente os mesmos atalhos podem ser usados para testar a implicação. Empregando uma noção semelhante de linhas ruins, você pode poupar uma enorme quantidade de trabalho.
 
\practiceproblems
\problempart
Usando atalhos, verifique se cada sentença é uma tautologia, uma contradição ou nenhuma das duas.
\begin{compactlist}
	\item $\enot B \eand B$ %contra
	\item $\enot D \eor D$ %taut
	\item $(A\eand B) \eor (B\eand A)$ %contingent
	\item $\enot[A \eif (B \eif A)]$ %contra
	\item $A \eiff [A \eif (B \eand \enot B)]$ %contra
	\item $\enot(A\eand B) \eiff A$ %contingent
	\item $A\eif(B\eor C)$ %contingent
	\item $(A \eand\enot A) \eif (B \eor C)$ %tautology
	\item $(B\eand D) \eiff [A \eiff(A \eor C)]$%contingent
\end{compactlist}

\chapter{Tabelas-verdade parciais}\label{s:PartialTruthTable}

Às vezes, não precisamos saber o que acontece em todas as linhas de uma tabela-verdade. Às vezes, basta uma ou duas linhas.

\paragraph{Tautologia}
Para mostrar que uma sentença é uma tautologia, precisamos mostrar que ela é verdadeira em toda valoração. Isto é, precisamos saber que ela resulta verdadeira em todas as linhas da tabela-verdade. Portanto, precisamos de uma tabela-verdade completa.

Para mostrar, no entanto, que uma sentença \emph{não} é uma tautologia, precisamos de apenas uma linha: uma linha na qual a sentença é falsa. Portanto, para mostrar que uma sentença não é uma tautologia, basta fornecer uma única valoração---uma única linha da tabela-verdade---que torne a sentença falsa.

Suponha que queiramos mostrar que a sentença `$(U \eand T) \eif (S \eand W)$' \emph{não} é uma tautologia. Montamos uma \define{tabela-verdade parcial}:
\begin{center}
\begin{tabular}{c c c c |d e e e e e f}
$S$&$T$&$U$&$W$&$(U$&\eand&$T)$&\eif    &$(S$&\eand&$W)$\\
\hline
   &   &   &   &    &   &    &\TTbf{F}&    &   &   
\end{tabular}
\end{center}
Deixamos espaço para apenas uma linha, em vez de 16, pois estamos procurando somente uma linha na qual a sentença seja falsa (daí também o `F').

O operador lógico principal da sentença é um condicional. Para que o condicional seja falso, o antecedente deve ser verdadeiro e o consequente deve ser falso. Portanto, preenchemos esses valores na tabela:
\begin{center}
\begin{tabular}{c c c c |d e e e e e f}
$S$&$T$&$U$&$W$&$(U$&\eand&$T)$&\eif    &$(S$&\eand&$W)$\\
\hline
   &   &   &   &    &  T  &    &\TTbf{F}&    &   F &   
\end{tabular}
\end{center}
Para que `$(U\eand T)$' seja verdadeira, tanto `$U$' quanto `$T$' devem ser verdadeiras.
\begin{center}
\begin{tabular}{c c c c|d e e e e e f}
$S$&$T$&$U$&$W$&$(U$&\eand&$T)$&\eif    &$(S$&\eand&$W)$\\
\hline
   & T & T &   &  T &  T  & T  &\TTbf{F}&    &   F &   
\end{tabular}
\end{center}
Agora só precisamos tornar `$(S\eand W)$' falsa. Para isso, precisamos tornar pelo menos uma de `$S$' e `$W$' falsa. Podemos tornar ambas, `$S$' e `$W$', falsas, se quisermos. O que importa é que a sentença inteira resulte falsa nessa linha. Tomando uma decisão arbitrária, terminamos a tabela desta maneira:
\begin{center}
\begin{tabular}{c c c c|d e e e e e f}
$S$&$T$&$U$&$W$&$(U$&\eand&$T)$&\eif    &$(S$&\eand&$W)$\\
\hline
 F & T & T & F &  T &  T  & T  &\TTbf{F}&  F &   F & F  
\end{tabular}
\end{center}
Agora temos uma tabela-verdade parcial, que mostra que `$(U \eand T) \eif (S \eand W)$' não é uma tautologia. Dito de outro modo, mostramos que existe uma valoração que torna `$(U \eand T) \eif (S \eand W)$' falsa, a saber, a valoração que torna `$S$' falsa, `$T$' verdadeira, `$U$' verdadeira e `$W$' falsa.

\paragraph{Contradições}
Mostrar que algo é uma contradição na LVF requer uma tabela-verdade completa: precisamos mostrar que não existe valoração que torne a sentença verdadeira; isto é, precisamos mostrar que a sentença é falsa em todas as linhas da tabela-verdade.

Para mostrar, no entanto, que algo \emph{não} é uma contradição, tudo o que precisamos fazer é encontrar uma valoração que torne a sentença verdadeira, e uma única linha de uma tabela-verdade será suficiente. Podemos ilustrar isso com o mesmo exemplo.
\begin{center}
\begin{tabular}{c c c c|d e e e e e f}
$S$&$T$&$U$&$W$&$(U$&\eand&$T)$&\eif    &$(S$&\eand&$W)$\\
\hline
  &  &  &  &   &   &   &\TTbf{T}&  &  &
\end{tabular}
\end{center}
Para tornar a sentença verdadeira, será suficiente garantir que o antecedente seja falso. Como o antecedente é uma conjunção, podemos simplesmente tornar um de seus conjuntivos falso. Tomando uma decisão arbitrária, vamos tornar `$U$' falso; então podemos atribuir qualquer valor de verdade que quisermos às outras letras sentenciais.
\begin{center}
\begin{tabular}{c c c c|d e e e e e f}
$S$&$T$&$U$&$W$&$(U$&\eand&$T)$&\eif    &$(S$&\eand&$W)$\\
\hline
 F & T & F & F &  F &  F  & T  &\TTbf{T}&  F &   F & F
\end{tabular}
\end{center}

\paragraph{Equivalência}
Para mostrar que duas sentenças são equivalentes, devemos mostrar que elas têm o mesmo valor de verdade em toda valoração. Portanto, isso requer uma tabela-verdade completa.

Para mostrar que duas sentenças \emph{não} são equivalentes, precisamos apenas mostrar que existe uma valoração na qual elas têm valores de verdade diferentes. Portanto, isso requer somente uma tabela-verdade parcial de uma linha: faça a tabela de modo que uma sentença seja verdadeira e a outra falsa.

\paragraph{Satisfatibilidade conjunta}
Para mostrar que algumas sentenças são conjuntamente satisfatíveis, devemos mostrar que existe uma valoração que torna todas as sentenças verdadeiras; portanto, isso requer apenas uma tabela-verdade parcial com uma única linha.

Para mostrar que algumas sentenças são conjuntamente insatisfatíveis, devemos mostrar que não existe valoração que torne todas as sentenças verdadeiras. Portanto, isso requer uma tabela-verdade completa: você deve mostrar que, em cada linha da tabela, pelo menos uma das sentenças é falsa.

\paragraph{Validade e implicação}
Para mostrar que um argumento é válido, devemos mostrar que não existe valoração que torne verdadeiras todas as premissas e falsa a conclusão. Portanto, isso requer uma tabela-verdade completa. (O mesmo vale para a implicação.)

Para mostrar que um argumento é \emph{inválido}, devemos mostrar que existe uma valoração que torna verdadeiras todas as premissas e falsa a conclusão. Portanto, isso requer apenas uma tabela-verdade parcial de uma linha na qual todas as premissas são verdadeiras e a conclusão é falsa. (O mesmo vale para uma falha de implicação.)

Esta tabela resume o que é necessário:

\begin{center}
\begin{tabular}{l l l}
%\cline{2-3}
 & \textbf{Sim} & \textbf{Não}\\
 \hline
%\cline{2-3}
tautologia? & completa & parcial de uma linha \\
contradição? &  completa & parcial de uma linha \\
%contingent? & two-line partial truth table & complete truth table\\
equivalente? & completa  & parcial de uma linha \\
satisfatível? & parcial de uma linha & completa \\
válido? & completa & parcial de uma linha \\
implicação? & completa & parcial de uma linha\\
\end{tabular}
\end{center}
\label{table.CompleteVsPartial}


\practiceproblems

\problempart
\label{pr.TT.equiv3}
Use tabelas-verdade completas ou parciais (conforme apropriado) para determinar se estes pares de sentenças são logicamente equivalentes:
\begin{compactlist}
\item $A$, $\enot A$ %No
\item $A$, $A \eor A$ %Yes
\item $A\eif A$, $A \eiff A$ %Yes
\item $A \eor \enot B$, $A\eif B$ %No
\item $A \eand \enot A$, $\enot B \eiff B$ %Yes
\item $\enot(A \eand B)$, $\enot A \eor \enot B$ %Yes
\item $\enot(A \eif B)$, $\enot A \eif \enot B$ %No
\item $(A \eif B)$, $(\enot B \eif \enot A)$ %Yes
\end{compactlist}

\problempart
\label{pr.TT.satisfiable4}
Use tabelas-verdade completas ou parciais (conforme apropriado) para determinar se estas sentenças são conjuntamente satisfatíveis ou conjuntamente insatisfatíveis:
\begin{compactlist}
\item $A \eand B$, $C\eif \enot B$, $C$ %unsatisfiable
\item $A\eif B$, $B\eif C$, $A$, $\enot C$ %unsatisfiable
\item $A \eor B$, $B\eor C$, $C\eif \enot A$ %satisfiable
\item $A$, $B$, $C$, $\enot D$, $\enot E$, $F$ %satisfiable
\item $A \eand (B \eor C)$, $\enot(A \eand C)$, $\enot(B \eand C)$ %satisfiable
\item $A \eif B$, $B \eif C$, $\enot(A \eif C)$ %unsatisfiable
\end{compactlist}

\problempart
\label{pr.TT.valid4}
Use tabelas-verdade completas ou parciais (conforme apropriado) para determinar se cada argumento é válido ou inválido:
\begin{compactlist}
\item $A\eor\bigl[A\eif(A\eiff A)\bigr] \therefore A$ %invalid
\item $A\eiff\enot(B\eiff A) \therefore A$ %invalid
\item $A\eif B, B \therefore A$ %invalid
\item $A\eor B, B\eor C, \enot B \therefore A \eand C$ %valid
\item $A\eiff B, B\eiff C \therefore A\eiff C$ %valid
\end{compactlist}

\problempart
\label{pr.TT.TTorC3}
Determine se cada sentença é uma tautologia, uma contradição ou uma sentença contingente. Justifique sua resposta com uma tabela-verdade completa ou parcial, conforme apropriado.

% truth tables in LaTeX generated by http://www.curtisbright.com/logic/. Be sure to give him a shout out.

\begin{compactlist}
\item  $A \eif \enot A$ \vspace{.5ex}							

%{\color{red}
%$
%\begin{array}{c|cccc}
%A&A&\eif&\enot&A\\\hline
%T&T&\mathbf{F}&F&T\\
%F&F&\mathbf{T}&T&F
%\end{array}
%$ 
%
%Contingent	 \vspace{6pt}
%}
%	T letter, 2 connectives
\item $A \eif (A \eand (A \eor B))$ \vspace{.5ex}	

%{\color{red}
%$
%\begin{array}{cc|ccc@{}ccc@{}ccc@{}c@{}c}
%A&B&A&\eif&(&A&\eand&(&A&\eor&B&)&)\\\hline
%T&T&T&\mathbf{T}&&T&T&&T&T&T&&\\
%T&F&T&\mathbf{T}&&T&T&&T&T&F&&\\
%F&T&F&\mathbf{T}&&F&F&&F&T&T&&\\
%F&F&F&\mathbf{T}&&F&F&&F&F&F&&
%\end{array}
%$
%
%Tautology \vspace{6pt}
%}
%			2 letters, 3 connectives

\item $(A \eif B) \eiff (B \eif A)$ 	\vspace{.5ex}				%
%
%{\color{red}
%$
%\begin{array}{cc|c@{}ccc@{}ccc@{}ccc@{}c}
%a&b&(&a&\rightarrow&b&)&\leftrightarrow&(&b&\rightarrow&a&)\\\hline
%T&T&&T&T&T&&\mathbf{T}&&T&T&T&\\
%T&F&&T&F&F&&\mathbf{F}&&F&T&T&\\
%F&T&&F&T&T&&\mathbf{F}&&T&F&F&\\
%F&F&&F&T&F&&\mathbf{T}&&F&T&F&
%\end{array}
%$
%
%Contingent \vspace{6pt}
%
%}
%		2 letters, 3 connectives

\item $A \eif \enot(A \eand (A \eor B)) $	\vspace{.5ex}	

%{\color{red}
%$
%\begin{array}{cc|cccc@{}ccc@{}ccc@{}c@{}c}
%a&b&a&\rightarrow&\enot&(&a&\eand&(&a&\eor&b&)&)\\\hline
%T&T&T&\mathbf{F}&F&&T&T&&T&T&T&&\\
%T&F&T&\mathbf{F}&F&&T&T&&T&T&F&&\\
%F&T&F&\mathbf{T}&T&&F&F&&F&T&T&&\\
%F&F&F&\mathbf{T}&T&&F&F&&F&F&F&&
%\end{array}
%$
%
%Contingent	\vspace{6pt}
%
%}
%
% 2 letters, 4 connectives

\item $\enot B \eif [(\enot A \eand A) \eor B]$\vspace{.5ex} 

%{\color{red}
%$
%\begin{array}{cc|cccc@{}c@{}cccc@{}ccc@{}c}
%a&b&\enot&b&\rightarrow&(&(&\enot&a&\eand&a&)&\eor&b&)\\\hline
%T&T&F&T&\mathbf{T}&&&F&T&F&T&&T&T&\\
%T&F&T&F&\mathbf{F}&&&F&T&F&T&&F&F&\\
%F&T&F&T&\mathbf{T}&&&T&F&F&F&&T&T&\\
%F&F&T&F&\mathbf{F}&&&T&F&F&F&&F&F&
%\end{array}
%$
%Contingent	 \vspace{6pt}
%
%}
%	2 letters, 5 connectives

\item $\enot(A \eor B) \eiff (\enot A \eand \enot B)$ \vspace{.5ex}

%{\color{red}
%$
%\begin{array}{cc|cc@{}ccc@{}ccc@{}ccccc@{}c}
%a&b&\enot&(&a&\eor&b&)&\leftrightarrow&(&\enot&a&\eand&\enot&b&)\\\hline
%T&T&F&&T&T&T&&\mathbf{T}&&F&T&F&F&T&\\
%T&F&F&&T&T&F&&\mathbf{T}&&F&T&F&T&F&\\
%F&T&F&&F&T&T&&\mathbf{T}&&T&F&F&F&T&\\
%F&F&T&&F&F&F&&\mathbf{T}&&T&F&T&T&F&
%\end{array}
%$
%
%Tautology \vspace{6pt}
%}
%2 letters, 6 connectives

\item $[(A \eand B) \eand C] \eif B$\vspace{.5ex}							
%
%{\color{red}
%$
%\begin{array}{ccc|c@{}c@{}ccc@{}ccc@{}ccc}
%a&b&c&(&(&a&\eand&b&)&\eand&c&)&\rightarrow&b\\\hline
%T&T&T&&&T&T&T&&T&T&&\mathbf{T}&T\\
%T&T&F&&&T&T&T&&F&F&&\mathbf{T}&T\\
%T&F&T&&&T&F&F&&F&T&&\mathbf{T}&F\\
%T&F&F&&&T&F&F&&F&F&&\mathbf{T}&F\\
%F&T&T&&&F&F&T&&F&T&&\mathbf{T}&T\\
%F&T&F&&&F&F&T&&F&F&&\mathbf{T}&T\\
%F&F&T&&&F&F&F&&F&T&&\mathbf{T}&F\\
%F&F&F&&&F&F&F&&F&F&&\mathbf{T}&F
%\end{array}
%$
%
%Tautology \vspace{6pt}
%}
%
%3 letters, 3 connectives

\item $\enot\bigl[(C\eor A) \eor B\bigr]$\vspace{.5ex} 						
%
%{\color{red}
%$
%\begin{array}{ccc|cc@{}c@{}ccc@{}ccc@{}c}
%a&b&c&\enot&(&(&c&\eor&a&)&\eor&b&)\\\hline
%T&T&T&\mathbf{F}&&&T&T&T&&T&T&\\
%T&T&F&\mathbf{F}&&&F&T&T&&T&T&\\
%T&F&T&\mathbf{F}&&&T&T&T&&T&F&\\
%T&F&F&\mathbf{F}&&&F&T&T&&T&F&\\
%F&T&T&\mathbf{F}&&&T&T&F&&T&T&\\
%F&T&F&\mathbf{F}&&&F&F&F&&T&T&\\
%F&F&T&\mathbf{F}&&&T&T&F&&T&F&\\
%F&F&F&\mathbf{T}&&&F&F&F&&F&F&
%\end{array}
%$
%
%Contingent \vspace{6pt}
%
%}
%	 	3 letters, 3 connectives

\item $\bigl[(A\eand B) \eand\enot(A\eand B)\bigr] \eand C$ \vspace{.5ex}	
%
%{\color{red}
%$
%\begin{array}{ccc|c@{}c@{}ccc@{}cccc@{}ccc@{}c@{}ccc}
%a&b&c&(&(&a&\eand&b&)&\eand&\enot&(&a&\eand&b&)&)&\eand&c\\\hline
%T&T&T&&&T&T&T&&F&F&&T&T&T&&&\mathbf{F}&T\\
%T&T&F&&&T&T&T&&F&F&&T&T&T&&&\mathbf{F}&F\\
%T&F&T&&&T&F&F&&F&T&&T&F&F&&&\mathbf{F}&T\\
%T&F&F&&&T&F&F&&F&T&&T&F&F&&&\mathbf{F}&F\\
%F&T&T&&&F&F&T&&F&T&&F&F&T&&&\mathbf{F}&T\\
%F&T&F&&&F&F&T&&F&T&&F&F&T&&&\mathbf{F}&F\\
%F&F&T&&&F&F&F&&F&T&&F&F&F&&&\mathbf{F}&T\\
%F&F&F&&&F&F&F&&F&T&&F&F&F&&&\mathbf{F}&F
%\end{array}
%$
%
%Contradiction \vspace{6pt}
%
%}
%
%% 	3 letters, 5 connectives
%
\item $(A \eand B) \eif[(A \eand C) \eor (B \eand D)]$ \vspace{.5ex}		
%
%{\color{red}
%$
%\begin{array}{cccc|c@{}c@{}ccc@{}c@{}ccc@{}c@{}ccc@{}ccc@{}ccc@{}c@{}c}
%a&b&c&d&(&(&a&\eand&b&)&)&\eif&(&(&a&\eand&c&)&\eor&(&b&\eand&d&)&)\\\hline
%T&T&T&T&&&T&T&T&&&\mathbf{T}&&&T&T&T&&T&&T&T&T&&\\
%T&T&F&F&&&T&T&T&&&\mathbf{F}&&&T&F&F&&F&&T&F&F&&\\
%\end{array}
%$
%
%Contingent \vspace{6pt}
%}
%
%	4 letters, 5 connectives
\end{compactlist}

\noindent\problempart
\label{pr.TT.TTorC4}
Determine se cada sentença é uma tautologia, uma contradição ou uma sentença contingente. Justifique sua resposta com uma tabela-verdade completa ou parcial, conforme apropriado.
\begin{compactlist}
\item  $\enot (A \eor A)$\vspace{.5ex}							%	Contradiction		1 letter, 2 connectives
\item $(A \eif B) \eor (B \eif A)$\vspace{.5ex}					%	Tautology			2 letters, 2 connectives
\item $[(A \eif B) \eif A] \eif A$\vspace{.5ex}					%	Tautology			2 letters, 3 connectives
\item $\enot[( A \eif B) \eor (B \eif A)]$\vspace{.5ex}			%	Contradiction		2 letters, 4 connectives
\item $(A \eand B) \eor (A \eor B)$\vspace{.5ex} 				%	Contingent		2 letters, 5 connectives
\item $\enot(A\eand B) \eiff A$\vspace{.5ex} 					%contingent			2 letters, 3 connectives
\item $A\eif(B\eor C)$\vspace{.5ex} 							%contingent			3 letters, 2 connectives
\item $(A \eand\enot A) \eif (B \eor C)$\vspace{.5ex} 			%tautology			3 letters, 4 connectives 
\item $(B\eand D) \eiff [A \eiff(A \eor C)]$\vspace{.5ex}			%contingent			4 letters, 4 connectives
\item $\enot[(A \eif B) \eor (C \eif D)]$\vspace{.5ex} 			% Contingent. 		4 letters, 4 connectives
\end{compactlist}



\noindent\problempart
Determine se cada um dos seguintes pares de sentenças é logicamente equivalente usando tabelas-verdade. Justifique sua resposta com uma tabela-verdade completa ou parcial, conforme apropriado.
\begin{compactlist}
\item $A$ e $A \eor A$
\item $A$ e $A \eand A$
\item $A \eor \enot B$ e $A\eif B$
\item $(A \eif B)$ e $(\enot B \eif \enot A)$
\item $\enot(A \eand B)$ e $\enot A \eor \enot B$
\item $ ((U \eif (X \eor X)) \eor U)$ e $\enot (X \eand (X \eand U))$
\item $ ((C \eand (N \eiff C)) \eiff C)$ e $(\enot \enot \enot N \eif C)$
\item $[(A \eor B) \eand C]$ e $[A \eor (B \eand C)]$
\item $((L \eand C) \eand I)$ e $L \eor C$
\end{compactlist}


\noindent\problempart
\label{pr.TT.satisfiable5}
Determine se cada coleção de sentenças é conjuntamente satisfatível ou conjuntamente insatisfatível. Justifique sua resposta com uma tabela-verdade completa ou parcial, conforme apropriado.
\begin{compactlist}
\item $A\eif A$, $\enot A \eif \enot A$, $A\eand A$, $A\eor A$ %satisfiable
\item $A \eif \enot A$, $\enot A \eif A$%unsatisfiable.
\item $A\eor B$, $A\eif C$, $B\eif C$ %satisfiable
\item $A \eor B$, $A \eif C$, $B \eif C$, $\enot C$ %	Insatisfiable
\item $B\eand(C\eor A)$, $A\eif B$, $\enot(B\eor C)$  %unsatisfiable
\item $(A \eiff B) \eif B$,  $B \eif \enot (A \eiff B)$, $A \eor B$  %	Consistent
\item $A\eiff(B\eor C)$, $C\eif \enot A$, $A\eif \enot B$ %satisfiable
\item  $A \eiff B$,  $\enot B \eor \enot A$,  $A \eif  B$ % Consistent
\item $A \eiff B$, $A \eif C$, $B \eif D$, $\enot(C \eor D)$ %consitent
\item $\enot (A \eand \enot B)$,  $B \eif \enot A$, $\enot B$   %Consistent
\end{compactlist}

\noindent\problempart Determine se cada argumento é válido ou inválido. Justifique sua resposta com uma tabela-verdade completa ou parcial, conforme apropriado.
\label{pr.TT.valid5} 
\begin{compactlist}
\item $A\eif(A\eand\enot A)\therefore \enot A$% valid
\item $A \eor B$, $A \eif B$, $B \eif A \therefore  A \eiff B$  % Valid
\item $A\eor(B\eif A)\therefore \enot A \eif \enot B$ %valid
\item $A \eor B$, $A \eif B$, $ B \eif A \therefore  A \eand B$ %valid
\item $(B\eand A)\eif C$, $(C\eand A)\eif B\therefore (C\eand B)\eif A$ % invalid
\item $\enot (\enot A \eor \enot B)$, $A \eif \enot C \therefore  A \eif (B \eif C)$ % invalid.
\item $A \eand (B \eif C)$, $\enot C \eand (\enot B \eif \enot A)\therefore C \eand \enot C$ % valid
\item $A \eand B$, $\enot A \eif \enot C$, $B \eif \enot D \therefore  A \eor B$ % Invalid
\item $A \eif B\therefore (A \eand B) \eor (\enot A \eand \enot B)$ % invalid
\item $\enot A \eif B$,$ \enot B \eif C $,$ \enot C \eif A \therefore  \enot A \eif (\enot B \eor \enot C) $% Invalid

\end{compactlist}

\noindent\problempart Determine se cada argumento é válido ou inválido. Justifique sua resposta com uma tabela-verdade completa ou parcial, conforme apropriado.
\label{pr.TT.valid6} 
\begin{compactlist}
\item $A\eiff\enot(B\eiff A)\therefore A$ % invalid
\item $A\eor B$, $B\eor C$, $\enot A\therefore B \eand C$ % invalid
\item $A \eif C$, $E \eif (D \eor B)$, $B \eif \enot D\therefore (A \eor C) \eor (B \eif (E \eand D))$ % invalid
\item $A \eor B$, $C \eif A$, $C \eif B\therefore A \eif (B \eif C)$ % invalid
\item $A \eif B$, $\enot B \eor A\therefore A \eiff B$ % valid
\end{compactlist}
